\chapter{Concordance with primary literature}
\label{chap:concordance}

\index{concordance|textbf}

This chapter is the constitution of the monograph.  When chapters
disagree, this chapter is right.  Its function is threefold: to state
the five main theorems with their precise status, to locate every
result against the prior literature, and to serve as the single
authoritative reference for what is proved, what is conditional,
and what remains conjectural.

\bigskip

\noindent\textbf{The five main theorems.}\quad
The algebraic engine rests on five theorems, all proved:

\begin{center}
\renewcommand{\arraystretch}{1.3}
\begin{tabular}{clp{7.5cm}l}
& \textbf{Theorem} & \textbf{Statement} & \textbf{Status} \\
\hline
\textbf{A} & Bar-cobar adjunction
  & $\barBch \dashv \Omegach$ on $\Ran(X)$,
    with Verdier intertwining
    \textup{(}Thm~\ref{thm:bar-cobar-isomorphism-main}\textup{)}
  & \ClaimStatusProvedHere \\
\textbf{B} & Bar-cobar inversion
  & $\Omegach(\barBch(\cA)) \xrightarrow{\sim} \cA$
    on the Koszul locus
    \textup{(}Thm~\ref{thm:higher-genus-inversion}\textup{)}
  & \ClaimStatusProvedHere \\
\textbf{C} & Complementarity
  & $Q_g(\cA) + Q_g(\cA^!) = H^*(\overline{\mathcal{M}}_g, Z(\cA))$
    \textup{(}Thm~\ref{thm:quantum-complementarity-main}\textup{)};
    upgraded to shifted-symplectic Lagrangian geometry
  & \ClaimStatusProvedHere \\
\textbf{D} & Modular characteristic
  & $\kappa(\cA)$ universal, additive, anti-symmetric,
    $\hat{A}$-genus generating function
    \textup{(}Thm~\ref{thm:genus-universality}\textup{)}
  & \ClaimStatusProvedHere \\
\textbf{H} & Chiral Hochschild
  & $\ChirHoch^*(\cA)$ polynomial, Koszul-functorial
    \textup{(}Thms~\ref{thm:w-algebra-hochschild},
    \ref{thm:virasoro-hochschild},
    \ref{thm:critical-level-cohomology}\textup{)}
  & \ClaimStatusProvedHere
\end{tabular}
\end{center}

\noindent
The unifying principle: the modular $L_\infty$ convolution algebra
$\mathfrak{g}^{\mathrm{mod}}_\cA$ carries a natural modular $L_\infty$
structure from the Feynman transform of the modular operad.  The dg~Lie
algebra of Definition~\ref{def:modular-convolution-dg-lie} is its strict
model $\Convstr$; the homotopy-invariant object is $\Definfmod(\cA)$
(Theorem~\ref{thm:modular-homotopy-convolution}).  The shadow Postnikov tower
$\Theta_\cA^{\le 2} \to \Theta_\cA^{\le 3} \to \Theta_\cA^{\le 4}
\to \cdots$ is the finite-order projection of the universal MC
element~$\Theta_\cA$; the proved content is the finite-order engine
(through arity~$4$), and the full unrestricted $\Theta_\cA$ exists
by the bar-intrinsic construction
(Theorem~\ref{thm:mc2-bar-intrinsic};
Theorem~\ref{thm:recursive-existence}).
Theorems~A--D+H are proved projections of the
scalar level~$\kappa$.

\begin{remark}[Scope]
This chapter contains the reader-facing mathematical content:
principal contributions, literature comparisons, the three-pillar
architecture, the three concentric rings, cross-volume bridges,
and the modular homotopy theory development.  The editorial
apparatus---the constitutional declaration, the conjecture
stratification and homotopy templates, the six research directions,
and the project-level status assessment---appears in
Chapter~\ref{chap:editorial-constitution}.
\end{remark}


\section{Principal contributions}\label{sec:principal-contributions}

\begin{enumerate}[label=\emph{(\arabic*)}]

\item \emph{Configuration space bar-cobar adjunction
  (Theorems~\ref{thm:bar-cobar-isomorphism-main}
  and~\ref{thm:higher-genus-inversion}).}
  The bar and cobar constructions for chiral algebras are realized
  via explicit integrals over Fulton--MacPherson compactifications
  $\overline{C}_n(X)$, with the bar differential given by residues
  at collision divisors and the cobar differential by inclusions
  of strata.  The adjunction is proved for all Koszul chiral algebras
  and all genera.  Theorem~A decomposes into three layers
  (Remark~\ref{rem:theorem-A-decomposition}):
  the fundamental theorem of chiral twisting morphisms
  ($\mathrm{A}_0$, Theorem~\ref{thm:fundamental-twisting-morphisms}),
  bar concentration
  ($\mathrm{A}_1$, Theorem~\ref{thm:bar-concentration}), and
  the Verdier-geometric duality
  ($\mathrm{A}_2$, Theorem~\ref{thm:bar-cobar-isomorphism-main}).
  This is the first complete proof of chiral Koszul
  duality via configuration space geometry.
  \emph{Quillen upgrade}
  (\S\ref{sec:concordance-three-pillars}): the bar-cobar adjunction
  $B_\kappa \dashv \Omega_\kappa$ is a \emph{Quillen equivalence}
  in the Vallette model structure on conilpotent
  chiral coalgebras~\cite[Theorem~2.1]{Val16}
  (Theorem~\ref{thm:quillen-equivalence-chiral}),
  and every $\mathrm{Ch}_\infty$-algebra
  rectifies to a strict chiral algebra
  (Corollary~\ref{cor:rectification-ch-infty}).
  The homotopy category of dg chiral algebras is thereby
  equivalent to $\mathrm{Ch}_\infty$-algebras with
  $\infty$-morphisms modulo homotopy.

\item \emph{Genus universality
  (Theorem~\ref{thm:genus-universality}).}
  For any Koszul chiral algebra~$\cA$, the genus-$g$ obstruction
  factors as $\mathrm{obs}_g(\cA) = \kappa(\cA) \cdot \lambda_g$,
  where $\kappa$ is an algebra-dependent constant and $\lambda_g$
  is the universal Faber--Pandharipande tautological class
  on~$\overline{\mathcal{M}}_g$.
  The genus-independent coefficient $\kappa(\cA)$ is determined
  by the genus-$1$ curvature.  This is the \emph{scalar} genus
  universality (Level~$0$ of Remark~\ref{rem:four-levels}): the
  first Chern class of the factorization homology bundle.
  The full genus-$g$ invariant---the cohomology
  $H^*(\barB^{(g)}(\cA), \Dg{g})$ and its variation over
  $\mathcal{M}_g$---is categorically richer
  (Theorem~\ref{thm:chain-modular-functor}).

\item \emph{Deformation-obstruction complementarity
  (Theorem~\ref{thm:quantum-complementarity-main}).}
  For a Koszul pair $(\cA, \cA^!)$, the genus-$g$ obstructions
  satisfy $\mathrm{obs}_g(\cA) + \mathrm{obs}_g(\cA^!) =
  [\text{characteristic class of } Z(\cA)]$ in $H^*(\overline{\mathcal{M}}_g)$.
  The Kodaira--Spencer map is constructed for all Koszul pairs
  and all genera (Theorem~\ref{thm:kodaira-spencer-chiral-complete}).
  This is the first formulation and proof of complementarity
  for chiral algebras.

\item \emph{$\Eone$-chiral Koszul duality
  (Theorem~\ref{thm:e1-chiral-koszul-duality}).}
  The bar-cobar adjunction extends to $\Eone$-chiral algebras
  (nonlocal vertex algebras with $R$-matrix braiding), covering
  Yangians and toroidal algebras.  This goes beyond the
  $\Einf$ (commutative/local) setting of BD and FG.

\item \emph{Explicit computations.}
  Complete Koszul duality computations for all standard families:
  Heisenberg, free fermion, $bc$-$\beta\gamma$, affine Kac--Moody
  ($\widehat{\fg}_k \leftrightarrow \widehat{\fg}_{-k-2h^\vee}$),
  Virasoro
  \textup{(}same-family complementarity partner
  $\mathrm{Vir}_{26-c}$ at M/S-level, with the
  $\mathcal{W}_\infty$ realization retained as project\textup{)},
  $\mathcal{W}_N$ algebras, Yangians.
  Genus-$g$ free energies computed through $g = 10$ with
  Bernoulli asymptotics.  Bar cohomology dimensions through
  degree~$10$.  These explicit results are new.

\item \emph{Derived oper space from the bar complex
  (Theorem~\ref{thm:oper-bar}).}
  At the critical level $k = -h^\vee$, the bar complex is uncurved
  and recovers the derived space of
  $\fg^\vee$-opers on the formal disc:
  $H^n(\barB(\widehat{\fg}_{-h^\vee}))
  \cong \Omega^n(\mathrm{Op}_{\fg^\vee}(D))$.
  The $H^0$ identification gives the Feigin--Frenkel
  center~\cite{Feigin-Frenkel}; the $H^1$ identification
  gives K\"ahler differentials of the oper space; and the
  full derived identification follows from the bar-Ext
  identification combined with the
  Frenkel--Teleman theorem~\cite{FT06}.
  This provides a vertex-algebraic construction of a
  central ingredient in the geometric Langlands programme,
  independent of the categorical approach of
  Gaitsgory et al.~\cite{GLC24} and
  Raskin~\cite{Raskin24}.

\item \emph{Chiral Hochschild cohomology
  (Theorems~\ref{thm:w-algebra-hochschild},
  \ref{thm:virasoro-hochschild},
  and~\ref{thm:critical-level-cohomology}).}
  For the principal $\Walg$-algebra $\Walg^k(\fg)$ at generic level:
  $\ChirHoch^*(\Walg^k(\fg)) \cong \mathbb{C}[\Theta_1, \ldots, \Theta_r]$
  with $\Theta_i$ in degree~$m_i + 1$ (exponents of~$\fg$).
  Rank~$1$ (Virasoro): strict period~$2$.
  Rank $r \geq 2$: quasi-period $\operatorname{lcm}(m_1+1, \ldots, m_r+1)$.
  At critical level $k = \critLevel$:
  exterior$\otimes$polynomial structure
  $\Lambda^*(P_1, \ldots, P_r) \otimes \mathbb{C}[\Theta_1, \ldots, \Theta_r]$
  (Fuks--Feigin--Tsygan).
  Functorial under Koszul duality via the Connes
  operator~$S$ (Theorem~\ref{thm:main-koszul-hoch},
  Corollary~\ref{cor:hochschild-cup-exchange}).
  This promotes the Hochschild results to a
  principal contribution~(Theorem~H).

\item \emph{Chain-level modular functor
  (Theorem~\ref{thm:chain-modular-functor}).}
  The Feynman transform of the Getzler--Kapranov modular operad
  assembles the genus-graded bar complex into a homotopy-coherent
  modular functor: to each surface a cochain complex, to each
  degeneration a chain map, to each consistency relation a chain
  homotopy.  Passing to~$H^0$ at integrable level recovers the
  space of conformal blocks.  This is the structure toward which
  Theorems~A--D build
  (Definition~\ref{def:modular-homotopy-theory-intro}):
  the scalar genus tower
  $F_g = \kappa(\cA) \cdot \lambda_g^{\mathrm{FP}}$
  is its first Chern class; the spectral discriminant~$\Delta_\cA$
  is its Chern character; the KZ/Hitchin connection on conformal
  blocks is the flat connection on its cohomology sheaf
  (Remark~\ref{rem:moduli-variation}).

\end{enumerate}

\begin{remark}[Geometric Langlands connection]%
\label{rem:oper-langlands-concordance}
\index{geometric Langlands!oper from bar complex}%
\index{oper!concordance summary}%
Item~\emph{(6)} is structurally significant.  The bar complex
$\barB(\widehat{\fg}_{-h^\vee})$ at critical level is a
\emph{derived enhancement} of the Feigin--Frenkel center: it
simultaneously encodes the commutative algebra of functions on
the oper space (in $H^0$), its cotangent complex (in $H^1$),
and higher derived structure (in $H^{\geq 2}$).  That the
bar construction---defined purely in terms of configuration
space integrals on curves---reproduces a key geometric
object in the Langlands programme
(Theorem~\ref{thm:oper-bar}) illustrates the reach of
chiral Koszul duality beyond abstract homotopy algebra.
\end{remark}

\begin{remark}[The twisting morphism as structural protagonist]
\label{rem:tau-protagonist}
\index{twisting morphism!as structural protagonist}
The canonical twisting morphism $\tau \colon \barB_X(\cA) \to \cA$
is the thread connecting the four main theorems
(Convention~\ref{conv:bar-coalgebra-identity}):
Theorem~A establishes $\tau$; Theorem~B inverts it;
Theorem~C identifies its genus-$g$ curvature; Theorem~D computes
its leading coefficient.
At higher genus, $\tau$ extends to a curved twisting cochain
$\tau_{\mathrm{mod}} = \tau + \Theta_\cA$
whose curvature $\kappa(\cA)\cdot\omega_g$ governs the entire
modular deformation tower
(Remark~\ref{rem:theta-modular-twisting}).
\end{remark}

\section{Relationship to Beilinson--Drinfeld}

\begin{tabular}{>{\raggedright\arraybackslash}p{0.38\textwidth}>{\raggedright\arraybackslash}p{0.52\textwidth}}
\textbf{Our Terminology} & \textbf{BD Terminology} \\
\hline
Chiral algebra & Chiral algebra (same) \\
$\Einf$-chiral algebra & Vertex algebra / Chiral algebra \\
$\Eone$-chiral algebra & No direct analog (generalization) \\
Chiral bracket $\mu$ & Chiral operation $\mu: j_* j^*(A \boxtimes A) \to \Delta_! A$ \\
Bar complex $\B(\cA)$ & Chiral chain complex (cf.\ Chevalley--Eilenberg for Lie case) \\
Factorization structure & Factorization algebra structure \\
Ran space $\Ran(X)$ & $\mathfrak{R}(X)$ in BD notation \\
\end{tabular}

\begin{remark}[Key extension]
Our $\Eone$-chiral algebras extend BD's chiral algebras by dropping skew-symmetry. BD's chiral algebras are our $\Einf$-chiral algebras.
\end{remark}

\section{Relationship to Francis--Gaitsgory}
\index{Francis--Gaitsgory}

\begin{tabular}{>{\raggedright\arraybackslash}p{0.38\textwidth}>{\raggedright\arraybackslash}p{0.52\textwidth}}
\textbf{Our Terminology} & \textbf{FG Terminology} \\
\hline
Chiral Koszul duality & Chiral Koszul duality (same) \\
$\chirCom$-$\chirLie$ duality & Main theorem of FG \\
$\chirAss$-$\chirAss$ duality & Not explicitly treated (implicit) \\
Pro-nilpotence & Pro-nilpotent tensor $\infty$-category \\
Bar-cobar equivalence & Koszul duality equivalence \\
\end{tabular}

\begin{remark}[Key extension]
FG establish $\chirCom$-$\chirLie$ duality. We show this is derived from the more fundamental $\chirAss$-$\chirAss$ self-duality via the deformation Pois → Ass.
\end{remark}

\section{Relationship to Gui--Li--Zeng}
\index{Gui--Li--Zeng}

\begin{tabular}{>{\raggedright\arraybackslash}p{0.38\textwidth}>{\raggedright\arraybackslash}p{0.52\textwidth}}
\textbf{Our Terminology} & \textbf{GLZ Terminology} \\
\hline
Quadratic chiral algebra & Quadratic chiral algebra (same) \\
Koszul dual $\cA^!$ & Quadratic dual $A^!$ \\
Bar complex & Not used (direct quadratic dual) \\
Non-quadratic duality & Not treated \\
$\Eone$-chiral algebras & Not treated \\
\end{tabular}

\begin{proposition}[GLZ as special case; \ClaimStatusProvedHere]\label{prop:glz-special-case}
The Gui--Li--Zeng quadratic duality is a special case of the $\chirAss$-duality of this monograph:
\[
\text{GLZ duality} = \chirAss\text{-duality} |_{\Einf\text{-chiral}} |_{\text{quadratic}}
\]
Restricting to $\Einf$-chiral algebras with quadratic presentations recovers their theory.
\end{proposition}

\begin{proof}
An $\Einf$-chiral algebra $\cA$ with quadratic presentation has a quadratic relation ideal $R \subset j_*j^*(\cA^{\boxtimes 2})$.  Its $\chirAss$ Koszul dual is determined by bar degrees $\leq 2$: the generators $V^*$ in $\bar{B}^1$ and the orthogonal complement of relations $R^\perp$ in $\bar{B}^2$ (the higher bar degrees, while nonzero, contribute no independent data for quadratic algebras).  The Koszul dual coalgebra $\bar{B}(\cA)$ is then determined by $R^\perp$, and the Koszul dual algebra $\cA^! = \bar{B}(\cA)^\vee$ is recovered by linear duality under the chiral pairing $j_*j^*(\cA^{\boxtimes 2}) \otimes j_*j^*((\cA^!)^{\boxtimes 2}) \to \omega_X$.  This is precisely the quadratic dual construction of GLZ.

For the $\Einf$ restriction: our $\chirAss$ bar complex uses the full associative chiral product $\mu\colon j_*j^*(\cA \boxtimes \cA) \to \Delta_!\cA$.  When $\cA$ is $\Einf$-chiral (i.e., the chiral bracket is skew-symmetric), the bar complex acquires a commutative coalgebra structure, and the Koszul dual is a Lie chiral algebra.  This recovers the FG framework as a special case (see Theorem~\ref{thm:fg-from-assch}).
\end{proof}

\begin{theorem}[FG duality from \texorpdfstring{$\chirAss$}{Ass-ch} self-duality; \ClaimStatusProvedHere]\label{thm:fg-from-assch}
The Francis--Gaitsgory $\chirCom$-$\chirLie$ duality is the associated graded of the $\chirAss$ self-duality under the PBW filtration.  Precisely: for a quadratic $\Einf$-chiral algebra $\cA$, the $\chirAss$ bar complex $\bar{B}_{\chirAss}(\cA)$ carries a filtration $F_\bullet$ (induced by the symmetrization degree of the $\Sigma$-action on $C_n(X)$) such that:
\begin{equation}\label{eq:fg-from-assch}
\mathrm{gr}^F \bar{B}_{\chirAss}(\cA) \;\cong\; \bar{B}_{\chirCom}(\mathrm{Sym}^{\mathrm{ch}}(\cA)) \;\cong\; \mathrm{CE}^{\chirLie}(\cA)
\end{equation}
and the associated spectral sequence degenerates at $E_1$ for Koszul algebras.
\end{theorem}

\begin{proof}
\emph{Step~1: PBW filtration.}
The chiral bar complex $\bar{B}^n_{\chirAss}(\cA)$ consists of sections on $\overline{C}_{n+1}(X)$ with coefficients in $\cA^{\boxtimes(n+1)} \otimes \Omega^n_{\log}$.  Define $F_p \bar{B}^n$ as the span of elements symmetric under at least $(n-p)$ transpositions of configuration space points.  This is the chiral analogue of the arity filtration on the classical associative bar complex (Loday--Vallette \cite{LV12}, \S9.3.2).

\emph{Step~2: Associated graded identification.}
On $\overline{C}_n(X)$, the $\Sigma_n$-representation decomposes as $j_*j^*(\cA^{\boxtimes n}) = \mathrm{Sym}^n_{\mathrm{ch}}(\cA) \oplus \bigwedge^2_{\mathrm{ch}}(\cA) \otimes \cdots$.  The PBW filtration separates these: $F_0/F_1 = \mathrm{Sym}^n_{\mathrm{ch}}(\cA)$ (the fully symmetric part) is the commutative chiral bar complex, and $F_1/F_2$ involves the Lie bracket (the antisymmetric part).  The associated graded $\mathrm{gr}^F \bar{B}_{\chirAss}$ is thus the Chevalley--Eilenberg complex for the Lie chiral algebra structure, which is $\bar{B}_{\chirCom}(\mathrm{Sym}^{\mathrm{ch}}(\cA))$ by the classical identification of the commutative bar complex with Chevalley--Eilenberg chains.

\emph{Step~3: Spectral sequence degeneration.}
For Koszul $\cA$, the spectral sequence $E_1^{p,q} = H^{p+q}(\mathrm{gr}^F_p \bar{B})$ degenerates at $E_1$.  To see this, note that the associated graded $\mathrm{gr}^F \bar{B}_{\chirAss}(\cA)$ is the Chevalley--Eilenberg complex (Step~2), and the Koszul property of $\cA$ as a chiral algebra implies that this commutative bar complex has cohomology concentrated on the diagonal $p + q = 0$ (i.e., $E_1^{p,q} = 0$ for $p + q \neq 0$).  Since $d_1\colon E_1^{p,q} \to E_1^{p+1,q}$ shifts the total degree $p+q$ by~$1$, it maps from the surviving diagonal to the zero region, hence $d_1 = 0$ and $E_2 = E_1$.  All higher differentials vanish for the same degree reason.  Therefore $H^*(\bar{B}_{\chirAss}(\cA)) \cong H^*(\mathrm{gr}^F \bar{B}_{\chirAss}(\cA))$, and the FG duality is recovered as the associated graded of our $\chirAss$ duality.

\emph{Step~4: Comparison with FG.}
The FG theorem states: for an $\Einf$-chiral algebra $\cA$, the chiral Koszul dual is a Lie chiral algebra $\cA^{!,\chirLie}$.  In the present framework: the Koszul dual coalgebra $\bar{B}_{\chirAss}(\cA)$ is an $\Eone$-chiral coalgebra, and the Koszul dual algebra $\cA^{!,\chirAss} = \bar{B}_{\chirAss}(\cA)^\vee$ is an $\Eone$-chiral algebra; its associated graded under PBW is $\cA^{!,\chirLie} = \bar{B}_{\chirCom}(\cA)^\vee$.  (The cobar $\Omega(\bar{B}_{\chirAss}(\cA))$ recovers~$\cA$ itself by bar-cobar inversion, not~$\cA^!$.)  The FG duality is recovered.
\end{proof}

\section{Relationship to Ayala--Francis}
\index{Ayala--Francis}

\begin{tabular}{>{\raggedright\arraybackslash}p{0.38\textwidth}>{\raggedright\arraybackslash}p{0.52\textwidth}}
\textbf{Our Terminology} & \textbf{AF Terminology} \\
\hline
Non-abelian Poincar\'{e} duality & Nonabelian Poincar\'{e} duality (same) \\
FM compactification $\overline{C}_n(X)$ & Disk stratification of $\mathrm{Ran}(X)$ \\
Verdier duality on $\overline{C}_n(X)$ & Poincar\'{e}/Koszul duality for $\Etwo$-algebras on oriented surfaces \\
Bar complex as factorization coalgebra & Factorization homology \\
\end{tabular}

\begin{remark}[Key comparison]\label{rem:ayala-francis-comparison}
The Ayala--Francis framework~\cite{AF15, AF20} establishes non-abelian
Poincar\'{e} duality as an equivalence between factorization homology and
factorization cohomology for framed $\mathsf{E}_n$-algebras on $n$-manifolds.
Our Theorem~A (Theorem~\ref{thm:bar-cobar-isomorphism-main}) is the
specialization to dimension~$n = 2$ (complex curves viewed as real
oriented surfaces) with the additional
structure of:
\begin{enumerate}[label=(\roman*)]
\item \emph{Holomorphic enrichment.}
  The AF framework works with topological $\mathsf{E}_n$-algebras and
  locally constant cosheaves.  We work with holomorphic chiral algebras
  ($\mathcal{D}$-modules on curves), where the Fulton--MacPherson
  compactification provides the correct holomorphic compactification
  (not just the topological one).  The bar differential involves
  residues of meromorphic forms, not just topological data.
\item \emph{Explicit formulas.}
  The AF duality is an $\infty$-categorical equivalence; our realization
  via configuration space integrals gives explicit chain-level formulas
  for the bar differential (residues at collision divisors) and the
  cobar differential (inclusions of FM strata).
\item \emph{Higher genus.}
  For genus $g \geq 1$ curves, the configuration space
  $C_n(\Sigma_g)$ acquires additional topology ($H^1(\Sigma_g)$),
  and the bar complex inherits new generators from the handle
  cycles.  The AF framework predicts this (factorization homology
  on non-simply-connected manifolds), but the explicit identification
  with moduli space integrals and the Hodge bundle (Theorem~C)
  is new.
\end{enumerate}
\end{remark}


\section{Relationship to Feigin--Frenkel}
\index{Feigin--Frenkel}

\begin{tabular}{>{\raggedright\arraybackslash}p{0.38\textwidth}>{\raggedright\arraybackslash}p{0.52\textwidth}}
\textbf{Our Terminology} & \textbf{FF Terminology} \\
\hline
KM Koszul dual $\widehat{\fg}_{-k-2h^\vee}$ & FF involution $k \mapsto -k - 2h^\vee$ \\
Critical level $k = -h^\vee$ & Critical level (same) \\
Curvature $\kappa \propto (k + h^\vee)$ & Center of $\widehat{\fg}_{-h^\vee}$ \\
DS reduction $\mathcal{W}^k(\fg)$ & Drinfeld--Sokolov reduction (same) \\
$c + c' = 2d$ & Not formulated \\
Genus-$g$ obstruction $\mathrm{obs}_g$ & Not formulated \\
\end{tabular}

\begin{remark}[Key comparison]\label{rem:feigin-frenkel-comparison}
The Feigin--Frenkel involution $k \mapsto -k - 2h^\vee$ is, in this
monograph's framework, a consequence of chiral Koszul duality via
Verdier duality on Fulton--MacPherson compactifications.  The
canonical treatment---including the geometric mechanism, the
critical-level recovery of the FF center, and the proof of
Theorem~\ref{thm:universal-kac-moody-koszul}---appears in
Chapter~\ref{chap:kac-moody}.

The extensions beyond FF are:
\begin{enumerate}[label=(\roman*)]
\item The complementarity formula $c(\widehat{\fg}_k) + c(\widehat{\fg}_{-k-2h^\vee}) = 2d$
  (Theorem~\ref{thm:central-charge-complementarity}), which shows the
  central charge sum is a root datum invariant, independent of the level.
\item The genus-$g$ obstruction formula $\mathrm{obs}_g = \kappa \cdot \lambda_g$
  (Theorem~\ref{thm:genus-universality}), which extends the FF involution
  to all genera via the Hodge bundle on $\overline{\mathcal{M}}_g$.
\item The Virasoro and $\mathcal{W}_N$ generalizations, where the
  Drinfeld--Sokolov parametrization introduces nonlinear level shifts
  (e.g., $c \mapsto 26 - c$ for Virasoro) that do not arise from a linear
  involution on levels.
\end{enumerate}
\end{remark}


\section{Relationship to Costello--Gwilliam}
\index{Costello--Gwilliam}

\begin{tabular}{>{\raggedright\arraybackslash}p{0.38\textwidth}>{\raggedright\arraybackslash}p{0.52\textwidth}}
\textbf{Our Terminology} & \textbf{CG Terminology} \\
\hline
Chiral algebra on $X$ & Holomorphic factorization algebra on $X$ \\
Bar complex $\B(\cA)$ & Factorization homology $\int_X \cA$ \\
Cobar construction $\Omega(C)$ & Factorization envelope \\
OPE residues & Observables in perturbative QFT \\
Envelope-shadow $\Thetaenv(R)$ & Shadow tower of $\Fact(R)$; \S\ref{sec:concordance-nishinaka-vicedo} \\
\end{tabular}

\begin{remark}[Key comparison]\label{rem:cg-comparison}
The Costello--Gwilliam factorization algebra framework~\cite{CG17}
relates to ours through the Beilinson--Drinfeld equivalence between
chiral algebras and factorization algebras on curves (BD, Chapter~3).
The precise comparison:
\begin{enumerate}[label=(\roman*)]
\item \emph{Factorization homology.}
  The CG factorization homology $\int_X \cA$ computes the same object
  as our chiral homology $H_*(\bar{B}^{\mathrm{ch}}(\cA))$.  Our explicit
  realization via configuration space integrals provides chain-level
  formulas that the CG framework treats abstractly.
\item \emph{Koszul duality.}
  The CG framework includes Koszul duality for factorization algebras
  (CG, Vol.~2, Chapter~5), specialized to holomorphic factorization
  algebras on curves.  Their Koszul duality, restricted to this setting,
  recovers our chiral Koszul duality.  Our contribution is the explicit
  identification of the Koszul dual for all standard families
  (Kac--Moody, Virasoro, $\mathcal{W}_N$, etc.) via the bar-cobar
  adjunction on Fulton--MacPherson spaces.
\item \emph{Quantum corrections.}
  The CG perturbative quantization framework produces $A_\infty$
  structures from Feynman diagram expansions (CG, Vol.~2, Chapter~3).
  Our curved $A_\infty$ structure on the bar complex is the
  chiral-algebraic counterpart: the curvature $m_0 = \kappa \cdot
  \mathbf{1}$ arises from the highest-order OPE pole (quartic for
  Virasoro, double for Heisenberg), not from a Feynman diagram sum.
  The genus universality theorem shows that these two perspectives
  produce the same genus-$g$ obstruction class.
\item \emph{BV-BRST.}
  The CG BV formalism produces the same BRST complex as our bar
  construction when applied to the chiral algebra of a 2d holomorphic
  field theory (Chapter~\ref{ch:bv-brst}).  The identification
  $\bar{B}(\cA) \simeq C^{\mathrm{BRST}}(\cA)$ connects the
  algebraic (bar complex) and physical (BRST cohomology) perspectives.
\item \emph{Monoidal bar-cobar.}
  Booth--Lazarev~\cite{BL24} construct monoidal model structures on
  categories of coalgebras, providing the homotopical infrastructure
  for the bar-cobar adjunction to be a Quillen equivalence in the
  monoidal setting.  This is realized by the fusion product preservation theorem
  (Theorem~\ref{thm:fusion-bar-cobar}): the bar construction
  is a monoidal functor from the tensor category of chiral
  algebra modules to the comodule category over the bar coalgebra.
\end{enumerate}
\end{remark}


\section{Relationship to Nishinaka--Vicedo}
\label{sec:concordance-nishinaka-vicedo}
\index{Nishinaka|textbf}
\index{Vicedo|textbf}
\index{factorization envelope!Lie conformal|textbf}
\index{enveloping vertex algebra|textbf}

Nishinaka~\cite{Nish26} constructs an amenably holomorphic
prefactorization algebra from a Lie conformal algebra~$R$ via a
\emph{factorization envelope}, and proves that the associated
vertex algebra is the enveloping vertex algebra~$\Uvert(R)$.
The construction explicitly recovers and generalizes the
Costello--Gwilliam Kac--Moody factorization algebra~\cite{CG17}
and Williams's Virasoro factorization algebra, and extends to
super Lie conformal algebras, producing the Neveu--Schwarz and
$\mathcal{N}=2$ vertex superalgebras.  Vicedo~\cite{Vic25}
gives a complementary unified factorization construction for
full universal enveloping vertex algebras, treating Kac--Moody,
Virasoro, and $\beta\gamma$ in one framework via unital
local Lie algebras.

In this form, factorization theory ceases to be a repository of
special examples and becomes a functorial machine:
\begin{equation}\label{eq:factorization-envelope-pipeline}
R \;\in\; \LCA
\quad\xrightarrow{\;\text{fact.\ env.}\;}\quad
\Fact(R)
\quad\xrightarrow{\;\text{assoc.\ VA}\;}\quad
\Uvert(R).
\end{equation}
Combined with the bar-cobar machinery of this monograph,
this pipeline acquires a nonlinear extension: the shadow
Postnikov tower becomes a functor on Lie conformal input,
not a case-by-case invariant of isolated vertex algebras.

\begin{definition}[Cyclically admissible Lie conformal algebra]
\label{def:cyclically-admissible}
\ClaimStatusProvedHere
\index{cyclically admissible|textbf}
\index{Lie conformal algebra!cyclically admissible}
A Lie conformal algebra~$L$ is \emph{cyclically admissible} if:
\begin{enumerate}[label=\textup{(\roman*)}]
\item it has a conformal-weight grading
  $L = \bigoplus_{h \geq 0} L_h$ with $\dim L_h < \infty$;
\item the descending filtration
  $F^m L := \bigoplus_{h \geq m} L_h$ is complete;
\item its OPE/$\lambda$-bracket has bounded pole order;
\item it carries an invariant residue pairing
  $\langle{-},{-}\rangle \colon L \otimes L \to \omega_X$
  compatible with translation and skew-symmetry.
\end{enumerate}
Conditions~(i)--(iii) are the Nishinaka admissibility
conditions ensuring that the genus-$0$ factorization
envelope~$\Fact_X(L)$ exists~\cite{Nish26}.
Condition~(iv) supplies the cyclic structure needed for the
cyclic deformation complex
$\Defcyc(\Uvert(L))$
(Definition~\ref{def:modular-cyclic-deformation-complex})
and modular completion.
\end{definition}

\begin{definition}[Envelope-shadow functor]
\label{def:envelope-shadow-functor}
\ClaimStatusProvedHere
\index{envelope-shadow functor|textbf}
\index{shadow Postnikov tower!envelope-shadow functor}
Let $R$ be a cyclically admissible Lie conformal algebra
(Definition~\ref{def:cyclically-admissible}).
Define the \emph{envelope-shadow functor} at arity~$r$ by
\begin{equation}\label{eq:envelope-shadow-functor}
\Thetaenv_{\leq r}(R)
\;:=\;
\Theta_{\leq r}\bigl(\Uvert(R)\bigr)
\;\in\;
\cA^{\mathrm{sh}}_r\bigl(\Uvert(R)\bigr),
\end{equation}
where $\Theta_{\leq r}$ is the arity-$r$ truncation of the
bar-intrinsic MC element
(Theorem~\textup{\ref{thm:mc2-bar-intrinsic}}).
The full envelope-shadow is
$\Thetaenv(R) := \varprojlim\, \Thetaenv_{\leq r}(R)$
(Theorem~\textup{\ref{thm:recursive-existence}}).
\end{definition}

\begin{definition}[Envelope-shadow complexity]
\label{def:envelope-shadow-complexity}
\ClaimStatusProvedHere
\index{envelope-shadow complexity|textbf}
\index{shadow depth!envelope-shadow complexity}
The \emph{envelope-shadow complexity} of~$R$ is
\begin{equation}\label{eq:envelope-shadow-complexity}
\chienv(R)
\;:=\;
r_{\max}\bigl(\Uvert(R)\bigr)
\;=\;
\sup\bigl\{r \ge 2 :
\cA^{\mathrm{sh}}_{r,0}\bigl(\Uvert(R)\bigr) \neq 0
\bigr\},
\end{equation}
the shadow depth
(Definition~\textup{\ref{def:shadow-depth-classification}})
of the enveloping vertex algebra.
\end{definition}

\begin{proposition}[Finite-jet rigidity]
\label{prop:finite-jet-rigidity}
\ClaimStatusProvedHere
\index{finite-jet rigidity|textbf}
\index{shadow Postnikov tower!finite-jet rigidity}
The arity-$r$ envelope-shadow $\Thetaenv_{\leq r}(R)$
depends only on the $(r{+}1)$-jet of the enveloping vertex
algebra --- that is, only on
$\Uvert(R) / F_{r+1}\,\Uvert(R)$,
where $F_\bullet$ is the PBW filtration.
\end{proposition}

\begin{proof}
The shadow tower at arity~$r$ is constructed from
$r$-vertex graph sums
(Definition~\ref{def:shadow-postnikov-tower}).
Each vertex carries a transferred operation
$\ell_{|v|}^{\mathrm{tr}}$ from the minimal model of
$\Defcyc(\cA)$, which at arity~$\leq r$ sees at most
$r{+}1$ OPE poles.  These are precisely the data carried
by $\Uvert(R)/F_{r+1}$.  Higher-jet corrections appear
only at arity $\geq r{+}1$.
\end{proof}

\begin{proposition}[Polynomial level dependence]
\label{prop:polynomial-level-dependence}
\ClaimStatusProvedHere
\index{shadow Postnikov tower!polynomial level dependence|textbf}
\index{central extension!polynomial dependence}
Let $R_t$ be a one-parameter family of Lie conformal algebras
with central-extension parameter~$t$ \textup{(}e.g., $t = k$
for affine Kac--Moody, $t = c$ for Virasoro\textup{)}.
Then $\Thetaenv_{\leq r}(R_t)$ is polynomial in~$t$ of
degree at most~$r{-}1$:
\begin{equation}\label{eq:polynomial-level}
\Thetaenv_{\leq r}(R_t)
\;=\;
\sum_{j=0}^{r-1} \theta_j^{(r)} \, t^j.
\end{equation}
In particular, $\kappa = \theta_1^{(2)} \, t$ is linear
\textup{(}Theorem~D\textup{)}, the cubic shadow
$\mathfrak{C} = \theta_0^{(3)} + \theta_1^{(3)} t +
\theta_2^{(3)} t^2$ is at most quadratic, and the quartic
resonance class~$\mathfrak{Q}$ is at most cubic in~$t$.
\end{proposition}

\begin{proof}
At arity~$r$, the shadow is a sum over stable graphs with
$\leq r$ vertices.  Each vertex contributes one factor of the
OPE structure constants, which are polynomial in~$t$ of
degree~$1$ (for affine/Virasoro type data) or degree~$0$
(for the structural part).  An $r$-vertex graph thus
contributes degree~$\leq r$ in~$t$.  The leading degree-$r$
coefficient factors through the scalar MC equation and is
absorbed into $\Theta_{\leq r-1}$ by the inductive
construction
(Construction~\ref{constr:obstruction-recursion}),
leaving degree~$\leq r{-}1$ for the genuinely new contribution
at arity~$r$.
\end{proof}

\begin{proposition}[Gaussian collapse for abelian input]
\label{prop:gaussian-collapse-abelian}
\ClaimStatusProvedHere
\index{Gaussian collapse|textbf}
\index{Heisenberg!Gaussian fixed point}
\index{envelope-shadow functor!abelian collapse}
Let $R$ be an abelian Lie conformal algebra
\textup{(}$[{-}_\lambda{-}]_R = 0$\textup{)}.  Then:
\begin{enumerate}[label=\textup{(\roman*)}]
\item $\Uvert(R)$ is a Heisenberg-type vertex algebra.
\item $\chienv(R) = 2$:
  the connected shadow tower terminates at arity~$2$.
\item $\Thetaenv(R) = \kappa(R) \cdot \eta \otimes \Lambda$
  \textup{(}scalar saturation\textup{)}.
\end{enumerate}
The Heisenberg algebra is the universal Gaussian fixed point
of the envelope-shadow functor: every abelian Lie conformal
input produces a class-$\mathbf{G}$ shadow.
\end{proposition}

\begin{proof}
The abelian bracket gives $\Uvert(R) = \mathrm{Sym}^{\mathrm{ch}}(R)$,
which is the Heisenberg algebra on the underlying $\DX$-module
with pairing determined by the cocycle of~$R$.
By the Heisenberg shadow computation
(\S\ref{sec:heisenberg-shadow-gaussianity}), the
shadow tower terminates at arity~$2$ with the scalar
curvature~$\kappa$.  Part~(iii) is scalar saturation
(Theorem~\ref{thm:algebraic-family-rigidity}).
\end{proof}

\begin{remark}[Envelope-shadow complexity of standard families]
\label{rem:envelope-shadow-complexity-standard}
The envelope-shadow complexity recovers the shadow depth
classification
(Definition~\ref{def:shadow-depth-classification}):
\begin{center}
\renewcommand{\arraystretch}{1.15}
\begin{tabular}{@{}llcl@{}}
\toprule
\textbf{$R$} & \textbf{$\Uvert(R)$}
  & $\chienv(R)$ & \textbf{Class} \\
\midrule
$R$ abelian & Heisenberg & $2$ & $\mathbf{G}$ \\
$\hat{\fg}_k$ current & $V_k(\fg)$ & $3$ & $\mathbf{L}$ \\
$\beta\gamma$ current & $\beta\gamma$ system & $4$ & $\mathbf{C}$ \\
Virasoro current & $\mathrm{Vir}_c$ & $\infty$ & $\mathbf{M}$ \\
\bottomrule
\end{tabular}
\end{center}
\end{remark}

\begin{remark}[From examples to machine]
\label{rem:from-examples-to-machine}
\index{factorization envelope!functorial completion}
The envelope-shadow functor upgrades the monograph's computational
programme from a census to a functor.  Classical treatments
construct each family's factorization algebra \emph{ad hoc}
(Heisenberg via free fields, Kac--Moody via currents, Virasoro
via the stress tensor).  The Nishinaka--Vicedo pipeline
\eqref{eq:factorization-envelope-pipeline} manufactures all of
them from a single input type (Lie conformal algebra), and the
bar-intrinsic construction
(Theorem~\ref{thm:mc2-bar-intrinsic}) then produces the shadow
tower $\Thetaenv(R)$ automatically.  The super extension
(Neveu--Schwarz, $\mathcal{N}=2$) suggests that the same
mechanism should organize much of the nonlinear chiral landscape.
\end{remark}

\begin{remark}[Conjectural frontier]
\label{rem:envelope-shadow-frontier}
\index{envelope-shadow functor!conjectural directions}
Five directions remain open:
\begin{enumerate}[label=\textup{(\roman*)},nosep]
\item \emph{DS-envelope descent.}
  If $R \twoheadrightarrow R_{\mathrm{DS}}$ is a
  Drinfeld--Sokolov reduction, does $\Thetaenv$ descend
  compatibly?  The expectation is
  $\Theta_{\mathcal{W}_N} = \mathrm{DS}(\Theta_{\widehat{\fg}})$:
  reduction should commute with the universal MC class,
  turning DS from an operation on individual examples into
  a functor on platonic packages
  (Construction~\ref{constr:platonic-package}).
  Reduction should happen \emph{before} shadow-taking:
  the shadow tower of $\mathcal{W}_N$ should be the
  BRST-reduced image of the affine tower.
\item \emph{Plumbing-Taylor theorem.}
  Identify the arity-$r$ shadow with the $r$-th Taylor
  coefficient of factorization homology
  $\int_{\Sigma_g} \Uvert(R)$ near a boundary divisor,
  using the plumbing-fixture expansion.
\item \emph{Lie-conformal complementarity.}
  Lift the complementarity theorem (Theorem~C) to a statement
  about $R$ and its Lie-conformal dual~$R^!$, giving
  $\Thetaenv(R) + \Thetaenv(R^!) = H^*(\overline{\cM}_g,
  Z(\Uvert(R)))$ at the Lie-conformal level.
\item \emph{Critical-level concentration.}
  At $k = -h^\vee$, the envelope
  $\Uvert(R_{-h^\vee})$ is uncurved; the shadow
  tower should concentrate on the factorizable
  Feigin--Frenkel center~\cite{CM25}.
\item \emph{Super shadow machine.}
  Extend $\Thetaenv$ to super Lie conformal algebras,
  computing the shadow tower for Neveu--Schwarz,
  $\mathcal{N}=2$, and the superconformal families.
\end{enumerate}
\end{remark}

\begin{construction}[The platonic package]
\label{constr:platonic-package}
\index{platonic package|textbf}
\index{factorization envelope!platonic package}
Fix a smooth complex curve~$X$ and a cyclically admissible Lie
conformal algebra~$L$
(Definition~\ref{def:cyclically-admissible}).
The \emph{platonic package} of~$L$ is the six-fold datum
\[
\Pi_X(L)
\;=\;
\bigl(
  \Fact_X(L),\;\;
  \barB_X(L),\;\;
  \Theta_L,\;\;
  \mathcal{L}_L,\;\;
  (V_L^{\mathrm{br}}, T_L^{\mathrm{br}}),\;\;
  \mathfrak{R}_4^{\mathrm{mod}}(L)
\bigr),
\]
where:
\begin{enumerate}[label=\textup{(\arabic*)}]
\item $\Fact_X(L)$ is the genus-$0$ factorization
  envelope~\cite{Nish26};
\item $\barB_X(L) := \barB(\Fact_X(L))$ is the reduced bar
  coalgebra;
\item $\Theta_L \in
  \MC(\Defcyc(L) \widehat\otimes \mathbb{G}_{\mathrm{mod}})$
  is the universal modular Maurer--Cartan class
  (Theorem~\ref{thm:mc2-bar-intrinsic});
\item $\mathcal{L}_L$ is the determinant line of the genus
  family;
\item $(V_L^{\mathrm{br}}, T_L^{\mathrm{br}})$ is the
  spectral branch object: a finite-rank complex with its
  genus-$1$ transport operator;
\item $\mathfrak{R}_4^{\mathrm{mod}}(L)$ is the quartic
  resonance class
  (Definition~\ref{def:nms-modular-quartic-resonance-class-restated}).
\end{enumerate}
The slogan is: \emph{envelope $+$ bar coalgebra $+$
universal MC class $=$ platonic form}.
Every invariant of the theory --- scalar anomaly, spectral
sheet, nonlinear resonance --- is a successive shadow of one
universal class~$\Theta_L$.
\end{construction}

\begin{conjecture}[The platonic adjunction]
\label{conj:platonic-adjunction}
\ClaimStatusConjectured
\index{platonic adjunction|textbf}
\index{factorization envelope!modular adjunction}
Let $\LCA_{\mathrm{cyc}}(X)$ denote the category of
cyclically admissible Lie conformal algebras on~$X$
(Definition~\ref{def:cyclically-admissible}),
and let $\mathsf{ModKoszul}(X)$ denote the
$\infty$-category of modular Koszul chiral/factorization
objects on~$X$.
Let
\[
\operatorname{Prim}^{\mathrm{mod}} \colon
\mathsf{ModKoszul}(X) \to \LCA_{\mathrm{cyc}}(X)
\]
be the primitive-current functor extracting the underlying
Lie conformal data together with its cyclic/modular
structure.  Then there exists a \emph{modular factorization
envelope} functor $U_X^{\mathrm{mod}}$ such that
\[
\operatorname{Hom}_{\mathsf{ModKoszul}(X)}
  \bigl(U_X^{\mathrm{mod}}(L),\, \cA\bigr)
\;\cong\;
\operatorname{Hom}_{\LCA_{\mathrm{cyc}}(X)}
  \bigl(L,\, \operatorname{Prim}^{\mathrm{mod}}(\cA)\bigr).
\]
At genus~$0$, this recovers the Nishinaka--Vicedo
envelope~\cite{Nish26,Vic25}.  In the full/non-chiral
face, it recovers Vicedo's construction~\cite{Vic25}.
The full modular envelope carries the platonic
package~$\Pi_X(L)$
(Construction~\ref{constr:platonic-package})
functorially.
\end{conjecture}

\begin{remark}[Execution programme]
\label{rem:envelope-execution-programme}
\index{factorization envelope!execution programme}
The factorization envelope programme decomposes into seven
stages:
\begin{enumerate}[label=\textup{(\arabic*)}]
\item Pin down the input category
  $\LCA_{\mathrm{cyc}}(X)$: Nishinaka-admissible plus
  cyclic pairing plus pronilpotent completion.
\item Build $\barB_X(L)$ functorially from the genus-$0$
  factorization envelope.
\item Construct $\Defcyc(L)$ as cyclic coderivations of
  $\barB_X(L)$, then transfer stable-graph operations
  to it.
\item Solve the Maurer--Cartan equation recursively by
  total complexity $2g - 2 + n + \text{bar weight}$,
  using the genus-$0$ envelope as the base point.
\item Extract the shadows in order:
  $\Theta_L \rightsquigarrow
  \kappa(L),\; \mathcal{L}_L,\;
  T_L^{\mathrm{br}},\; \Delta_L,\;
  \mathfrak{R}_4^{\mathrm{mod}}(L)$.
\item Compute the archetypes:
  Heisenberg, $\widehat{\mathfrak{sl}}_2$,
  $\beta\gamma$, Virasoro,
  principal~$\mathcal{W}_N$.
\item Pass from the $E_\infty$/local case to the ordered
  $\Eone$ face, where the same package acquires an
  $R$-matrix/Yang--Baxter refinement.
\end{enumerate}
Steps~(1)--(2) are supplied by the current literature.
Steps~(3)--(5) are the content of the proved core
(Theorem~\ref{thm:mc2-bar-intrinsic};
Definition~\ref{def:shadow-algebra};
Theorem~\ref{thm:recursive-existence}).
Steps~(6)--(7) are the constructive and spectral layers
developed throughout Parts~I--III.
\end{remark}


\section{Relationship to Loday--Vallette}

\begin{tabular}{>{\raggedright\arraybackslash}p{0.38\textwidth}>{\raggedright\arraybackslash}p{0.52\textwidth}}
\textbf{Our Terminology} & \textbf{LV Terminology} \\
\hline
Operad & Operad (same) \\
Koszul operad & Koszul operad (same) \\
Bar construction $\B$ & Bar construction $\B$ (same) \\
Cobar construction $\Omega$ & Cobar construction $\Omega$ (same) \\
Twisting morphism & Twisting morphism (same) \\
Maurer--Cartan equation & Maurer--Cartan equation (same) \\
$\Ass^! \cong \Ass$ & $\Ass^! \cong \Ass \otimes \mathrm{sgn}$ (Thm 7.1.1) \\
$\Com^! \cong \Lie$ & $\Com^! \cong \Lie\{1\}$ (Thm 7.2.1) \\
\end{tabular}
\begin{remark}[Convention for Koszul dual isomorphisms]
In Loday--Vallette \cite{LV12}, \S7.1--7.2, the Koszul dual includes the sign representation: $\Ass^! = \Ass \otimes \mathrm{sgn}$ and $\Com^! = \Lie\{1\}$ (with operadic suspension). In this manuscript, we absorb the sign twist into the bar desuspension: our bar construction uses the shifted generators $s^{-1}\mathcal{A}$ whose sign already accounts for the $\mathrm{sgn}$-twist. Concretely, our convention $\Ass^! \cong \Ass$ means that the dual cooperad is isomorphic to $\Ass$ as a \emph{graded} cooperad after this shift. All bar-cobar computations in the text use the desuspended convention consistently, so no sign errors propagate. See Appendix~\ref{app:signs} for the explicit translation between conventions.
\end{remark}

\begin{remark}[Chiral enhancement]
We lift the entire LV framework to the chiral setting, replacing vector spaces with D-modules and tensor products with chiral tensor products.
\end{remark}

\begin{remark}[Overview]\label{rem:concordance-synthesis}
The common thread running through the comparisons above is that each
framework---Beilinson--Drinfeld's chiral algebras, Francis--Gaitsgory's
chiral Koszul duality, Costello--Gwilliam's factorization algebras,
Nishinaka--Vicedo's factorization envelopes,
Loday--Vallette's operadic Koszul theory---captures one projection of a
single underlying phenomenon: \emph{the duality between commutative and
Lie structures, mediated by configuration spaces}.  The contribution of
the present monograph is the explicit chain-level realization of this
duality via configuration space integrals, which makes all projections
simultaneously computable.  The bar complex $\bar{B}^{\mathrm{ch}}(\cA)$
is the meeting point: it is BD's chiral homology, FG's Koszul dual
coalgebra, CG's factorization homology on the Ran space, NV's
factorization envelope of Lie conformal input, and LV's bar
construction---all incarnated as a single cochain complex with an
explicitly computable differential.
At the homotopy level, $\Convinf$ unifies all these frameworks:
each named framework is a strictification or truncation of the
modular $L_\infty$-deformation object $\Definfmod(\cA)$.
\end{remark}

\section{Relationship to Malikov--Schechtman}
\label{sec:concordance-malikov-schechtman}
\index{Malikov--Schechtman|textbf}
\index{homotopy chiral algebra!concordance}

Malikov--Schechtman~\cite{MS24} introduce \emph{homotopy chiral algebras}
(HCA): chiral algebras up to higher homotopies, with the Moore complex of
a cosimplicial Lie algebra in a pseudo-tensor category carrying $L_\infty$
operations controlled by the Eilenberg--Zilber operad
$\mathcal{Y}(n) = \mathcal{Z}(n) \otimes_k \mathrm{Lie}(n)$.  Their main
theorem (Theorem~3.1): the \v{C}ech complex of a sheaf of chiral algebras
is an HCA\@.  In our language (Definition~\ref{def:cech-convolution}):
the HCA operations assemble into a single MC element
$\Phi_\cA \in \operatorname{MC}(\mathfrak{g}^{\check{C}}_{\cA})$
satisfying $D\Phi + \tfrac{1}{2}[\Phi,\Phi] = 0$.

\begin{tabular}{>{\raggedright\arraybackslash}p{0.38\textwidth}%
>{\raggedright\arraybackslash}p{0.52\textwidth}}
\textbf{Our framework} & \textbf{MS framework} \\
\hline
Convolution algebra $\mathfrak{g}^{\mathrm{bar}}_\cA$
  & Convolution algebra $\mathfrak{g}^{\check{C}}_\cA$ \\
MC element $\Theta_\cA = (m_1, m_2, \ldots)$
  & MC element $\Phi_\cA = (F_1, F_2, \ldots)$ \\
$A_\infty$ operations $\{m_n\}$
  & HCA operations $\{F_n\}$ \\
Ran space / Fulton--MacPherson
  & Zariski site / affine cover \\
Associahedra / moduli $\overline{\mathcal{M}}_{0,n}$
  & Eilenberg--Zilber operad $\mathcal{Y}(n)$ \\
Modular operad FCom (genus $g \geq 1$)
  & [Not treated; genus-$0$ only] \\
$m_3$ from $\FM_3(\mathbb{C})$ integration
  & $F_3$ from triple intersections \\
$d^2 = 0$ on $\overline{\mathcal{M}}_{0,n}$
  & Cosimplicial identities \\
\end{tabular}

\begin{remark}[Complementary resolutions]
\label{rem:ms-complementary}
The bar and \v{C}ech complexes resolve the same cohomology with
compatible MC structures.  The comparison
(Proposition~\ref{prop:cech-bar-intertwining}) is an $L_\infty$
morphism $\Phi^*$ sending $\Phi_\cA$ to $\Theta_\cA$; it is
generically \emph{non-strict} because the bar side encodes
configuration space topology ($\FM_k$ strata) while the \v{C}ech
side encodes cover topology (intersections).  The MS framework is
purely algebraic; ours works at all genera via FCom and the shadow
Postnikov tower.  The curvature $d^2 = \kappa \cdot \omega_g$ and
finite-order shadows have no \v{C}ech counterpart in~\cite{MS24}.
\end{remark}

\begin{remark}[Genus-$0$ strictness and the asymmetry]
\label{rem:ms-genus0-agreement}
On $\mathbb{P}^1$ with a two-element cover, the \v{C}ech HCA is
\emph{strict for every chiral algebra}
(Proposition~\ref{prop:two-element-strict}): the MC element is
$\Phi = (d_M, F_2, 0, 0, \ldots)$ because $\check{C}^p = 0$ for
$p \geq 2$ kills all higher operations.  The bar $A_\infty$
structure, by contrast, has $m_3 \neq 0$ for non-abelian algebras
(e.g., $\widehat{\mathfrak{sl}}_2$) because the triple collision
stratum of $\FM_3(\mathbb{C})$ contributes regardless of the cover.
The comparison $\Phi^*$ creates $m_3$ from $F_2$ via its higher
components $\Phi_2^*$.  Genuine \v{C}ech non-strictness ($F_3 \neq 0$)
appears at genus~$g \geq 1$, where covers require $\geq 3$ opens and
triple intersections contribute
(Example~\ref{ex:cech-hca-genus1}).
\end{remark}

\section{Three-pillar foundational architecture}
\label{sec:concordance-three-pillars}
\index{three-pillar architecture|textbf}
\index{homotopy chiral algebra!foundational role}
\index{convolution!sL-infinity@$sL_\infty$-algebra!foundational role}
\index{logarithmic Fulton--MacPherson!foundational role}

Three recent papers force upgrades at three different levels of the
monograph's foundations.  We record the resulting \emph{three-pillar
architecture} as a constitutional amendment.

The three pillars are not independent imports.  They form a triangle
whose center is the master MC equation
$D\cdot\Theta_\cA + \tfrac{1}{2}[\Theta_\cA,\Theta_\cA] = 0$.
Each pillar provides one face of~$\Theta_\cA$: Pillar~A gives its
\emph{local algebraic nature} (a homotopy chiral algebra on the
curve), Pillar~B gives its \emph{deformation-theoretic home} (the
convolution $sL_\infty$-algebra that controls all twisting
morphisms), and Pillar~C gives its \emph{global geometric skeleton}
(the log-FM compactification that provides the collision geometry).
Every theorem in the monograph can be read through all three lenses
simultaneously; the three-pillar architecture is a single
refoundation, not three separate additions.

\begin{remark}[The three pillars]
\label{rem:three-pillars}
\index{three-pillar architecture!summary}
The primitive objects of the theory are:
\begin{enumerate}[label=\textup{(\Alph*)}]
\item \textbf{Local algebraic object} (Pillar~A).  The primitive local
  object is a \emph{homotopy chiral algebra}
  ($\mathrm{Ch}_\infty$-algebra) in the sense of
  Malikov--Schechtman~\cite{MS24}, not a strict chiral algebra.  A strict
  chiral algebra appears only after rectification or on cohomology.
  The chain-level story is
  $\mathrm{Ch}_\infty \to \text{strict chiral algebra on } H^\bullet
  \to \text{PVA/coisson shadow}$.
  The PVA descent from $A_\infty$ chiral algebra to $(-1)$-shifted
  Poisson vertex algebra on cohomology is proved in Volume~II
  (Theorem~\ref{thm:cohomology_PVA} in Section~\ref{sec:PVA_descent}):
  commutativity, skew-symmetry, Jacobi, and Leibniz (D2--D5) are
  established by exchange-cylinder and three-face Stokes arguments
  on~$\FM_3(\mathbb{C})$; higher-operation vanishing (D6) is proved
  via H3(a) factorization plus topological contractibility of
  $\operatorname{Conf}_k^{<}(\mathbb{R})$.
\item \textbf{Deformation-control object} (Pillar~B).  The universal
  deformation machine is a \emph{filtered convolution $sL_\infty$-algebra}
  $\operatorname{hom}_\alpha(\cC, \cA)$ relative to a twisting morphism
  $\alpha$, in the sense of Robert-Nicoud--Wierstra~\cite{RNW19} and
  Vallette~\cite{Val16}.  The deformation space is the Maurer--Cartan
  $\infty$-groupoid $\mathrm{MC}_\bullet(\operatorname{hom}_\alpha)$,
  not a naive MC set modulo gauge.
  The dg~Lie algebra $\mathfrak{g}_\cA^{\mathrm{mod}}$ of
  Definition~\ref{def:modular-convolution-dg-lie} is the strict model of
  the homotopy-invariant $\Definfmod(\cA)$
  (Theorem~\ref{thm:modular-homotopy-convolution}).
\item \textbf{Global collision geometry} (Pillar~C).  The natural
  geometry for punctured curves, nodal degenerations, and modular sewing
  is \emph{logarithmic Fulton--MacPherson} compactification
  $\overline{\operatorname{FM}}_n(X|D)$ on simple normal crossings pairs
  $(X, D)$, in the sense of Mok~\cite{Mok25}.  Ordinary FM
  compactification is the special case $D = \emptyset$.  The
  tropicalization of $\overline{\operatorname{FM}}_n(X|D)$ is controlled
  by planted-forest combinatorics
  (Definition~\ref{def:planted-forest-coefficient-algebra}).
\end{enumerate}
\end{remark}

\begin{remark}[The triangle]
\label{rem:three-pillar-triangle}
\index{three-pillar architecture!triangle}
The three pillars form a coherent triangle.  The bar complex
$\barB(\cA)$ is simultaneously:
\begin{itemize}
\item a $\mathrm{Ch}_\infty$-algebra (Pillar~A, via cosimplicial structure
  and the Eilenberg--Zilber operad $\mathcal{Y}(n)$);
\item an algebra in
  $\operatorname{hom}_\alpha(\cC^{\textup{!`}}_{\mathrm{ch}},
  \cP^{\mathrm{ch}})$ (Pillar~B, via the chiral operadic twisting
  morphism);
\item a sheaf on $\overline{\operatorname{FM}}_n(C|D)$ for a curve~$C$
  with divisor~$D$ (Pillar~C, via log-smooth compactification of
  operator insertions).
\end{itemize}
The master Maurer--Cartan equation
$D \cdot \Theta_\cA + \tfrac{1}{2}[\Theta_\cA, \Theta_\cA] = 0$
(proved at all arities, Theorem~\ref{thm:mc2-bar-intrinsic})
lives at the intersection of all three structures.  Each face of the
triangle links two pillars:
\begin{enumerate}[label=\textup{(\roman*)}]
\item A+B: the $\infty$-category of chiral algebras
  (\cite[Theorem~3.8]{Val16} gives simplicial structure; \cite{MS24}
  gives operations);
\item B+C: the geometric convolution algebra (planted forests of
  Mok~\cite{Mok25} are the combinatorial data controlling the
  convolution bracket on $\mathfrak{g}_\cA^{\mathrm{mod}}$);
\item A+C: the geometric origin of higher homotopies (secondary Borcherds
  operations $j'_{(p,q,r)}$ are residues on FM boundary strata; log-FM
  degeneration organizes these across genera).
\end{enumerate}
\end{remark}


\begin{summary}[Concrete chain-level rewrite]
\label{sum:three-pillars-chain-level}
The abstract triangle already has a concrete strict chart.  Choose a
\v{C}ech homotopy chiral model
$A^{\mathrm{ch}}_\infty=\check C^\bullet(\mathfrak U;\cA)$ from
Malikov--Schechtman, choose Mok's log-FM chain cooperad
$\cC^{\log FM}_{\mathrm{mod}}=\{C_\bullet(\operatorname{FM}_n(X\mathbin{|}D))\}$,
and form the one-slot convolution object
\[
\Convinf\!\bigl(
\cC^{\log FM}_{\mathrm{mod}},
\operatorname{End}_{\mathrm{Ch}_\infty}(A^{\mathrm{ch}}_\infty)
\bigr).
\]
Its strict chart is the bar-coderivation dg~Lie algebra, its Feynman
expansion is the stable-graph formula for~$\Theta_\cA$, and its
differential splits into the internal, sewing, planted-forest, and loop
terms of~\eqref{eq:D-five-components}.  Weight~$1$ in this strict
chart is already concrete: the three boundary points of
$\overline{\cM}_{0,4}$ give the secondary Borcherds/Jacobiator homotopy,
while the irreducible boundary of $\overline{\cM}_{1,1}$ gives the
genus-one BV loop operator.  Weight~$2$ is the first rigid
planted-forest layer, so it is the first genuinely logarithmic stratum.
Thus the three pillars already determine the concrete chain-level
programme: modular bar functor, convolution $L_\infty$-object, graph
expansion, and clutching calculus.  The detailed formulas are collected
in \S\ref{subsec:three-preprints-chain-level} and
Section~\ref{sec:chain-level-modular-convolution-refoundation}: in fixed
bidegree the Maurer--Cartan equation becomes a three-source recursion
(separating gluing, one-edge loops, rigid planted forests), while the
end-state of the programme is that $\kappa$, $\Delta_\cA$,
$\mathfrak{C}$, $\mathfrak{Q}$, and the higher shadows are all
characteristic projections of the same twisted tangent complex at
$\Theta_\cA$.
\end{summary}

\begin{summary}[Platonic end-state: modular Chern--Weil factorization]
\label{sum:platonic-end-state}
The three papers make four concrete structures canonical:
\[
A^{\mathrm{ch}}_\infty,
\qquad
\Delta^{\log}_\Gamma,
\qquad
\ell_\Gamma,
\qquad
T^{\mathrm{mod}}_{\Theta_\cA}.
\]
These are, respectively, the derived chiral input, the log-FM boundary
cooperad, the graphwise Taylor coefficients, and the twisted tangent
complex at the universal modular Maurer--Cartan point.  The end-state
formula is
\[
\Theta_\cA
=
\sum_\Gamma \frac{W^{\log FM}_\Gamma}{|\Aut(\Gamma)|}\,\Phi^\cA_\Gamma,
\qquad
 d_{\Theta_\cA}(x)
 =
 \sum_{n,g}\frac{\hbar^g}{n!}\,\ell_{n+1}^{(g)}(\Theta_\cA^{\otimes n},x),
\]
and the familiar invariants are extracted by
\[
\kappa(\cA)=\operatorname{tr}(\Theta_\cA)_{2,0},
\qquad
\Delta_\cA(t)=\operatorname{sdet}\!\bigl(1-t\,H(d_{\Theta_\cA}|_{\mathrm{red}})\bigr),
\qquad
\mathfrak R^{\mathrm{mod}}_{4,g,n}(\cA)=\operatorname{pr}_{4,g,n}(\Theta_\cA).
\]
Thus the scalar, spectral, cubic, quartic, and clutching packages are
not parallel themes but successive characteristic shadows of a single
modular convolution object.
\end{summary}


\begin{summary}[Derived modular atlas]
\label{sum:derived-modular-atlas}
The three references provide an actual atlas on the modular homotopy problem.
The Malikov--Schechtman chart is the operadic one
\[
\tau_{\le 0}Y(n)=\tau_{\le 0}\bigl(Z(n)\otimes \operatorname{Lie}(n)\bigr)
\curvearrowright
\check C^\bullet(\mathfrak U;\cA),
\]
the Robert-Nicoud--Wierstra chart is the Taylor-coefficient formula
\[
\ell_k(f_1,\dots,f_k)(x)=\gamma_A(\alpha\circ 1_D)(f_1\otimes\cdots\otimes f_k)\Delta_D^k(x),
\]
and the Mok chart is the rigid boundary subdivision
\[
\kappa\colon \Pi_n(\Delta)(\operatorname{cone}(\rho))\to \prod_v \Pi_{I_v}(\Sigma_v),
\qquad
\FM_n(W/B)(\rho)\to \prod_v \FM_{I_v}(Y_v\mathbin{|}D_v).
\]
Together these formulas explain why the modular bar filtration is simultaneously an arity
filtration, a cooperadic-coproduct filtration, and a codimension filtration on logarithmic
boundary strata.
\end{summary}

\begin{remark}[Key identifications]
\label{rem:three-pillar-identifications}
\index{three-pillar architecture!identifications}
The integration yields eleven identification theorems (six proved via
the cited sources; one structural obstruction from
Robert-Nicoud--Wierstra):
\begin{enumerate}[label=\textup{(\arabic*)}]
\item $\mathfrak{g}_\cA^{\mathrm{mod}} =
  \operatorname{hom}_\alpha(\cC^{\textup{!`}}_{\mathrm{ch}},
  \cP^{\mathrm{ch}})$ (convolution = master object;
  Theorem~\ref{thm:convolution-master-identification}).
\item $\Theta_\cA \in \mathrm{MC}(\operatorname{hom}_\alpha)
  \cong \mathrm{Tw}_\alpha$ (MC element = twisting morphism;
  Corollary~\ref{cor:theta-twisting-morphism}).
\item $\mathbb{G}_{\mathrm{pf}} =
  \mathrm{Trop}(\overline{\operatorname{FM}}_n(C|D))$
  (planted forests = tropicalization;
  Theorem~\ref{thm:planted-forest-tropicalization}).
\item $\barB(\cA)$ is a $\mathrm{Ch}_\infty$-algebra
  (Theorem~\ref{thm:cech-hca}).
\item $B_\kappa \dashv \Omega_\kappa$ is a Quillen equivalence
  (Theorem~\ref{thm:quillen-equivalence-chiral}).
\item The shadow algebra $\cA^{\mathrm{sh}}$ is a homotopy invariant
  (Theorem~\ref{thm:shadow-homotopy-invariance}).
\item The one-slot obstruction constrains the MC3 categorical lift:
  $\operatorname{hom}_\alpha$ accepts $\infty$-morphisms in either slot
  separately but not both simultaneously
  (\cite[Theorem~6.6]{RNW19}).
\item The secondary Borcherds operations $F_n$ of \cite{MS24} at
  arity $n \geq 3$ are identified with the shadow tower obstruction
  classes $o_n(\cA)$
  (Proposition~\ref{prop:borcherds-shadow-identification}):
  $F_3 = o_3$ (cubic shadow), $F_4 = o_4$ (quartic resonance),
  $F_n = o_n$ (higher obstructions).  The identification is
  structural at genus~$0$; at genus $g \geq 1$ the bar side
  acquires modular corrections with no \v{C}ech counterpart
  in~\cite{MS24}.
\item The quartic clutching law
  (Theorem~\ref{thm:nms-clutching-law-modular-resonance}) has a
  log-geometric derivation via Mok's degeneration formula
  \cite[Corollary~5.3.4]{Mok25}: the three terms (contact,
  additive, tree) biject with the three classes of rigid
  combinatorial types at arity~$4$
  (Construction~\ref{constr:arity4-degeneration}).
\item The log clutching conjecture is proved:
  Theorem~\ref{conj:log-clutching-degeneration} (formerly a
  conjecture) is now established by \cite[Corollary~5.3.4]{Mok25}.
\item The ambient $D^2 = 0$
  (Theorem~\ref{thm:ambient-d-squared-zero}, formerly
  Theorem~\ref{conj:differential-square-zero}) is now
  \textbf{proved}: the five-component differential on
  $\mathfrak{g}^{\mathrm{amb}}_{\cA}$ is the transport of
  $\partial$ on $\operatorname{FM}_n^{\mathrm{log}}(X|D)$
  through the Hom functor, and $\partial^2 = 0$ by Mok's
  normal-crossings result \cite[Theorem~3.3.1(1)]{Mok25}.
\end{enumerate}
\emph{Scope note.}  Identifications~(1)--(6) are stated at all
arities; however, the fullness of the all-arity unification
$\Theta_\cA \in \mathrm{MC}(\operatorname{hom}_\alpha)$ at all orders
is MC2 (proved by the bar-intrinsic construction,
Theorem~\ref{thm:mc2-bar-intrinsic}); finite-order projections
$\Theta_\cA^{\leq r}$ are its truncations.  Item~(7) is an external structural
fact independent of arity.
\end{remark}

\begin{remark}[Design constraint: failure of two-sided bifunctoriality]
\label{rem:no-bifunctor}
\index{bifunctoriality!failure of two-sided}
\index{one-slot obstruction}
Robert-Nicoud--Wierstra~\cite[Section~6]{RNW19} prove that although
the convolution functor $\operatorname{hom}_\alpha(-,-)$ extends to
$\infty_\alpha$-morphisms in either slot separately, there is
\emph{no} honest bifunctor accepting $\infty$-morphisms in both slots
simultaneously.  This is a structural constraint, not a technical
annoyance.  It means:
\begin{itemize}
\item Bar-cobar extends to $\infty$-morphisms of chiral algebras (slot~2)
  and to $\infty$-morphisms of chiral coalgebras (slot~1), but not both
  simultaneously.
\item The MC3 categorical lift
  (Conjecture~\ref{conj:mc3-arbitrary-type}) must proceed one slot at a
  time.  The prefundamental CG closure
  (Proposition~\ref{prop:prefundamental-clebsch-gordan}) works at the
  character level ($K_0$) precisely because characters live in one slot;
  the categorical lift requires both.
\item Any claim of two-sided functoriality for bulk/boundary/line
  constructions must be replaced by a homotopy-coherent correspondence
  formalism.
\end{itemize}
\end{remark}

\begin{remark}[Three-pillar architecture and the MC open problems]
\label{rem:three-pillar-mc-unification}
\index{three-pillar architecture!MC open problems}
\index{MC open problems!three-pillar unification}
The three-pillar architecture determines the MC resolution strategy:
\begin{itemize}
\item \emph{MC1 and MC2}:  Pillar~B realizes every algebra as an MC
  element in the modular convolution $sL_\infty$-algebra; the
  finite-order projections $\Theta_\cA^{\leq r}$ (the proved shadow
  Postnikov tower) are the constructive content.
\item \emph{MC3}:  MC3 is proved in type~$A$ (Theorem~\ref{thm:mc3-type-a-resolution}); beyond type~$A$, the categorical lift requires two-sided functoriality
  in $\operatorname{hom}_\alpha(-,-)$, which fails by the one-slot
  obstruction (Remark~\ref{rem:no-bifunctor}).  The workaround:
  character-level ($K_0$) generation via Pillar~B in one slot;
  categorical splitting awaits either rectification from Pillar~A
  or geometric resolution from Pillar~C.
\item \emph{MC4}:  The formal completion theory is \textbf{proved} by
  the strong completion-tower theorem
  (Theorem~\ref{thm:completed-bar-cobar-strong}): the finite-stage
  bar-cobar duality passes to inverse limits automatically once the
  filtration is strong (arity cutoff, Lemma~\ref{lem:arity-cutoff}).
  The completion closure $\CompCl(\Fft)$ carries a quasi-inverse
  bar-cobar equivalence on the homotopy category
  (Corollary~\ref{cor:completion-closure-equivalence}), stable under
  MC~twisting (Theorem~\ref{thm:mc-twisting-closure}), with completed
  twisting representability
  (Theorem~\ref{thm:completed-twisting-representability}).
  Splits into MC4$^+$ (positive towers) and MC4$^0$ (resonant).
  MC4$^+$ is \textbf{unconditionally solved} by weightwise
  stabilization: for $\mathcal{W}_{1+\infty}$, affine Yangians, and
  positive RTT towers, coefficient stabilization
  (Theorem~\ref{thm:coefficient-stability-criterion}) reduces
  completion to finite truncations.  MC4$^0$ concerns algebras with
  finite resonance rank---Virasoro, non-quadratic
  $\mathcal{W}_N$---and is reduced to a finite-dimensional problem
  by Theorem~\ref{thm:resonance-filtered-bar-cobar}.  The uniform
  PBW bridge (Theorem~\ref{thm:uniform-pbw-bridge}) connects MC1
  to MC4: uniform PBW degeneration plus quotient stabilization
  implies completed bar-cobar duality.
  The remaining example-specific task is coefficient stabilization on
  finite windows and H-level target identification, not a formal
  completion theorem.
\item \emph{MC5}:  The BV/BRST identification at genus~$\leq 1$ is
  proved; genus~$\geq 2$ gains a new approach via Pillar~C: Mok's
  log-smooth degeneration provides an inductive decomposition of the
  genus-$g$ bar complex through rigid special-fibre components.
\end{itemize}
\end{remark}

\section{The three concentric rings}
\label{sec:three-rings}
\index{three concentric rings|textbf}
\index{nonlinear characteristic layer|textbf}

The monograph has crossed a threshold.  What began as a theory of
bar-cobar duality on curves has grown into a structure with three
concentric rings.

\subsection*{Ring~1: The proved modular Koszul core}
Theorems~A/B/C/D/H, MC1/MC2, DK-0 through DK-3 on the
evaluation-generated core.  This is the spine.  Nothing in
the extension appendices touches it except to cross-reference
from it.

\subsection*{Ring~2: The nonlinear characteristic layer}
Before the extension appendices, the modular characteristic was
essentially scalar: the single number~$\kappa(\cA)$ and its
$\lambda_g$ tower, plus the spectral
discriminant~$\Delta_\cA(x)$.  Now there is a hierarchy:
\[
\kappa(\cA)
\;\longrightarrow\;
\Delta_\cA
\;\longrightarrow\;
[\mathfrak{C}_\cA],\;
[\mathfrak{Q}_\cA]
\;\longrightarrow\;
\Lambda_P(\mathfrak{Q}_\cA)
\;\longrightarrow\;
\Theta_\cA.
\]
\begin{itemize}
\item The \emph{nonlinear modular shadows} chapter
  (Chapter~\ref{app:nonlinear-modular-shadows}) makes the
  cubic and quartic jets of the complementarity
  potential~$S_\cA$ into precise, computable invariants with
  clutching laws and a quartic resonance class.
\item The \emph{all-arity master equation} extends the shadow
  tower to arbitrary arity~$r$: the primitive archetypes terminate
  (Gaussian at~$2$, Lie/tree at~$3$, contact at~$4$), while
  Virasoro is the first genuinely infinite tower
  (Theorem~\ref{thm:nms-finite-termination}).
\item The \emph{genus loop operator}~$\Lambda_P$ implements
  non-separating clutching: the genus-$1$ Hessian correction
  $\delta H^{(1)}=\Lambda_P(\mathfrak Q^{(0)})$ is the first
  tensorial modular invariant beyond the $\hat{A}$-genus
  (Theorem~\ref{thm:nms-beyond-ahat}).  For the Virasoro,
  $\delta H^{(1)}_{\mathrm{Vir}}
  =120/[c^2(5c+22)]\cdot x^2$.
\item The \emph{branch-line reductions} appendix
  (Appendix~\ref{app:branch-line-reductions}) extracts exact
  finite-dimensional quotients: the universal genus-$3$
  constant~$7$, genus-$2$ transparency, and the spectral
  cumulant hierarchy.
\item The \emph{ambient complementarity upgrade}
  (\S\ref{sec:ambient-complementarity-lagrangian}) upgrades
  Theorem~C from additive splitting to shifted-symplectic
  Lagrangian geometry: one ambient deformation problem with two
  Lagrangian maps, a complementarity potential~$S_\cA$ whose
  Taylor jets are the nonlinear modular invariants, a derived
  critical locus governing simultaneous self-dual deformations,
  and a fake-complementarity criterion detecting when the dual
  Lagrangian is linear
  (Proposition~\ref{prop:fake-complementarity-criterion}).
\item The \emph{modular bar--Hamiltonian}
  (\S\ref{sec:modular-bar-hamiltonian}) identifies the carrier
  algebra~$\mathfrak{g}^{\mathrm{amb}}_\cA$ with its five-piece
  differential (internal $+$ kernel twist $+$ stable-edge sewing
  $+$ planted-forest correction $+$ genus loop).  The universal
  class~$\Theta_\cA$ is a Maurer--Cartan element in this algebra
  (Definition~\ref{def:modular-bar-hamiltonian}), and the recursive
  extension tower provides the constructive mechanism for building
  it order by order
  (Construction~\ref{constr:obstruction-recursion}).
\item The \emph{modular Chern--Weil transform}
  (\S\ref{sec:modular-chern-weil-transform}) is a single map from
  $\operatorname{MC}(\mathfrak{g}^{\mathrm{amb}}_\cA)$ to the full
  invariant hierarchy: scalar classes, spectral classes, resonance
  classes, and line-kernel data are all projections of
  one Maurer--Cartan element.  The complementarity
  Legendrian~$\Lambda_\cA = j^1\mathfrak{S}_\cA$, the
  complementarity Stokes groupoid, and the phase
  index~$\nu_{\mathrm{nl}}(\cA,u)$ are geometric avatars of the
  same master element.
\item The \emph{quartic resonance class}
  $\mathfrak{R}_{4,g,n}^{\mathrm{mod}}(\cA)$
  (Theorem~\ref{thm:nms-clutching-law-modular-resonance}) is the
  first Picard-valued nonlinear characteristic of the theory, with
  clutching law receiving a cubic tree correction
  $\mathfrak{T}_{3,\xi}(\cA)$ on separating boundaries.
\item The \emph{graph-completed Feynman transform package}
  (\S\ref{sec:nms-feynman-transform-programme}) identifies the
  carrier object~$\mathfrak{F}_{\mathrm{mod}}(\cA)$ in which the
  full universal class should live, the full
  \emph{resonance tower}~$\mathfrak{R}^{\mathrm{mod}}_{r}$ for all
  arities~$r \ge 4$, and the \emph{nonlinear
  spectral-branch stack}~$\mathfrak{B}^{\mathrm{nl}}_\cA$.
\item The \emph{shadow archetypes in Part~II}: each standard
  family chapter now computes its nonlinear shadow explicitly.
  Heisenberg is exactly Gaussian
  (\S\ref{sec:heisenberg-shadow-gaussianity}); affine
  $\widehat{\mathfrak{sl}}_2$ is Lie/tree with critical-level
  emergence (\S\ref{sec:affine-cubic-shadow}); $\beta\gamma$ is
  the contact/quartic primitive with rank-one abelian rigidity
  (\S\ref{sec:betagamma-quartic-birth}); Virasoro and
  $\mathcal{W}_N$ are the first mixed cubic--quartic families with
  infinite shadow towers (\S\ref{sec:mixed-cubic-quartic-shadows}).
  The Virasoro quartic contact coefficient
  $Q^{\mathrm{contact}}_{\mathrm{Vir}} = 10/[c(5c+22)]$
  and the genus-$1$ loop ratio
  $\rho^{(1)}_{\mathrm{Vir}} = 240/[c^3(5c+22)]$ are the first
  non-scalar Ring~2 invariants extracted in closed form.
\item The \emph{envelope-shadow functor}
  (\S\ref{sec:concordance-nishinaka-vicedo}) upgrades the shadow
  tower from a case-by-case invariant to a functor on Lie conformal
  input $R \mapsto \Thetaenv(R) := \Theta(\Uvert(R))$, using the
  Nishinaka--Vicedo factorization-envelope pipeline.  The
  envelope-shadow complexity $\chienv(R)$ recovers the
  G/L/C/M classification; the arity-$r$ shadow depends only on
  the $(r{+}1)$-jet (Proposition~\ref{prop:finite-jet-rigidity})
  and is polynomial of degree~$\leq r{-}1$ in the
  central-extension parameter
  (Proposition~\ref{prop:polynomial-level-dependence}).
\item The \emph{universal Maurer--Cartan principle}
  (\S\ref{sec:nms-universal-mc-principle}) establishes the modular
  convolution algebra~$\mathfrak{g}_\cA^{\mathrm{mod}}$ as a modular
  $L_\infty$ algebra (not merely dg~Lie), with the shadow tower
  $\kappa \to \Delta \to \mathfrak{C} \to \mathfrak{Q} \to \Theta$
  arising as evaluation of the universal MC element~$\Theta_\cA$
  on specific moduli cycles.  The five-component differential
  $D = d_{A_\infty} + d_{\mathrm{sew}} + \hbar\Delta + d_r$
  is automatic from the single MC equation.  The dg~Lie model
  is the strict shadow of the $L_\infty$ structure; bar-cobar
  preserves quasi-isos because it is a quantum $L_\infty$ functor.
\item The \emph{non-renormalization theorem}
  (Theorem~\ref{thm:non-renormalization-tree}, ProvedElsewhere,
  Dimofte--Niu--Py) validates all tree-level computations: for
  quasi-linear theories (all HT-twisted 3d $\mathcal{N}=2$), the
  $A_\infty$ operations and the $r$-matrix are exact at tree level.
  Loop corrections enter only through the genus-raising BV
  operator~$\Delta$.
\item The \emph{polarized scattering package}
  (\S\ref{sec:nms-polarized-scattering-programme}) replaces the
  one-colour Maurer--Cartan problem by a bipartite machine:
  a polarized modular graph
  algebra~$\mathcal{G}^{\pm}_{\mathrm{mod}}(\cA)$ with
  same-colour vanishing, a \emph{resonance scattering
  diagram}~$\mathfrak{Scat}(\cA)$, a categorified
  \emph{complementarity schober}~$\mathfrak{Sch}(\cA)$, and a
  \emph{modular wavefunction}~$\Psi_\cA$ satisfying the quantum
  master equation.  Four conjectural theorems (bipartite
  vanishing, quartic lift, resonance wall-crossing, quantization)
  organize the forward attack.
\end{itemize}

\subsection*{Ring~3: The physics-facing extension problems}
Three axes, each with its own appendix chain:
\begin{enumerate}
\item \emph{The $\mathcal{W}$-algebra axis.}
  Completed Koszulity is ubiquitous (every affine
  $\mathcal{W}$-algebra at generic level); strict Koszulity is
  exceptional.  The subregular family
  $\mathcal{W}_n^{(2)}$ exhibits unbounded canonical homotopy
  arity.
\item \emph{The Yangian/RTT axis and the Bridge Theorem.}
  Three layers: local RTT kernel, completed envelope, boundary
  quotient.  The shift belongs to the boundary quotient, not
  the local kernel.  The rank-one test case produces the
  quantized Kleinian surface.
  The \emph{Bridge Theorem project}
  (Remark~\ref{rem:bridge-theorem-programme};
  Theorem~\ref{thm:bridge-criterion})
  is now the central remaining target: the full triple
  equivalence
  $\mathrm{Fact}^{\Eone}_{\mathrm{ord}} \simeq
  \mathrm{Mod}^{\mathrm{comp}}(\Ydg_\cA) \simeq
  \mathrm{Rep}^{\mathrm{spec}}(\QGspec(R_\cA))^{\mathrm{op}}$
  decomposes into four sub-targets~B1--B4 with a proved criterion
  theorem (B1$+$B2$+$B4 $\Rightarrow$ full bridge).
\item \emph{The holographic/celestial axis.}
  Protected holographic calculus (reconstruction, transposition,
  Weyl sewing, metaplectic anomaly) built directly from the
  proved core.
  Volume~II develops the bulk--boundary--line Koszul triangle
  $A_{\mathrm{bulk}} \simeq Z_{\mathrm{der}}(\text{boundary})
  \simeq \mathrm{HH}^\bullet(\cA^!)$.
  The \emph{holographic modular Koszul datum}
  $\mathcal{H}(T) = (\cA, \cA^!, \mathcal{C}, r(z), \Theta_\cA,
  \nabla^{\mathrm{hol}})$
  (Definition~\ref{def:holographic-modular-koszul-datum})
  packages the full HT holographic system into a single modular
  MC problem: the dg-shifted Yangian $r(z)$ is the binary
  genus-zero collision residue of~$\Theta_\cA$, and the 2026
  exact sphere-correlator matching~\cite{GZ26} is the genus-zero
  matrix-element shadow of this datum.  Five theorem targets
  decompose the programme: boundary--defect realization
  (Conjecture~\ref{conj:boundary-defect-realization}),
  Yangian-shadow
  (Conjecture~\ref{conj:yangian-shadow-theorem}),
  sphere reconstruction
  (Conjecture~\ref{conj:sphere-reconstruction}),
  quartic resonance obstruction
  (Conjecture~\ref{conj:quartic-resonance-obstruction}),
  and singular-fiber descent
  (Conjecture~\ref{conj:singular-fiber-descent}).
  The factorisation quantum-group envelope~\cite{Latyntsev23}
  supplies the categorical home beyond the meromorphic line category,
  and the HT deformation-quantization functor
  (Conjecture~\ref{conj:ht-deformation-quantization})
  bridges classical PVA data to the quantum modular Koszul package.
  The landscape of worked physics examples now includes:
  free chiral vortex lines with explicit $r$-matrix
  (\S\ref{sec:betagamma-vortex-lines}),
  affine PVA $\to$ Chern--Simons action
  (\S\ref{sec:affine-ht-chern-simons}),
  SQED--XYZ mirror symmetry as bar-cobar equivalence,
  non-renormalization of tree-level operations
  (\S\ref{sec:non-renormalization-tree}),
  superpotential $A_\infty$ truncation
  (\S\ref{sec:superpotential-ainfty-truncation}),
  Virasoro phase space as Teichm\"uller geometry
  (\S\ref{subsec:virasoro-teichmuller-phase-space}),
  level-rank duality as Koszul duality,
  monopole operators via affine Grassmannian modifications,
  Costello's M2-brane double-loop model
  (\S\ref{sec:m2-brane-double-loop}),
  SuperVirasoro $\to$ 3d supergravity,
  SQCD boundary algebra,
  Witten diagrams in twisted holography,
  $N=4$ global symmetry matching,
  the deformed conifold,
  and $\mathcal{W}_N \to$ higher-spin gauge theory.
\end{enumerate}

\subsection*{The unifying principle}
The single deepest structural insight across all three rings:
\emph{one object, one equation, everything else is a corollary.}
The modular $L_\infty$ convolution algebra~$\mathfrak{g}_\cA^{\mathrm{mod}}$
is the ambient home.  An algebra structure on~$\cA$ IS a Maurer--Cartan
element~$\Theta_\cA \in \mathrm{MC}(\mathfrak{g}_\cA^{\mathrm{mod}})$.
The shadow tower, the genus tower, the graph sums, the non-renormalization
theorems, and the shadow archetypes are all automatic from this single
definition.  The dg~Lie model is the strict computational avatar;
the $L_\infty$ structure is the reason bar-cobar is a Quillen equivalence.
Filtration controls the extension tower:
In Ring~2, the complementarity potential has a Taylor filtration
whose jets are the nonlinear shadows.  In Ring~3, the MC4 splitting
(Remark~\ref{rem:mc4-positive-vs-resonant}) separates
the completion problem:
positive towers ($\mathcal{W}_{1+\infty}$, Yangians) are
solved by weight stabilization
(Theorem~\ref{thm:stabilized-completion-positive});
resonant towers (Virasoro, non-quadratic $\mathcal{W}_N$) require the
resonance-filtered construction
(Theorem~\ref{thm:resonance-filtered-bar-cobar}),
with the same-family shadow $\mathrm{Vir}_{26-c}$ reinterpreted as
the depth-zero resonance shadow
(Remark~\ref{rem:virasoro-resonance-model}).
In Ring~3 (holography),
the composite filtration separates the wedge from the nonlinear
obstruction tower.

The five main theorems were about \emph{existence of duality}.
The three rings are about \emph{structure of the invariants
that duality produces}: from ``the machine exists'' to ``here
is what the machine computes, and here is how to read the
output.''


%% ====================================================================
%%  INTRINSIC KOSZULNESS: THE CHARACTERIZATION PROGRAMME
%% ====================================================================
\section{Intrinsic Koszulness: the characterization programme}
\label{sec:concordance-koszulness-programme}
\index{Koszul property!characterization programme|textbf}
\index{chiral Koszulness!intrinsic characterization|textbf}

The PBW criterion (Theorem~\ref{thm:pbw-koszulness-criterion})
detects chiral Koszulness for every standard family but does
not characterize it intrinsically.  This section is the
constitutional index for the intrinsic characterization
programme.  Full statements and proofs live in the source
chapters cited below.

%% ------------------------------------------------------------------
\subsection{The meta-theorem: twelve characterizations}
\label{subsec:concordance-koszulness-meta-theorem}
\index{Koszul property!equivalences|textbf}

Let $\cA$ be an augmented chiral algebra with PBW
filtration~$F_\bullet$.  The full meta-theorem is
Theorem~\ref{thm:koszul-equivalences-meta} in
Chapter~\ref{chap:koszul-pairs}.

\medskip\noindent
\textbf{Tier~1: proved equivalences}
($\textup{(i)} \Leftrightarrow \textup{(ii)}
\Leftrightarrow \textup{(iii)}
\Leftrightarrow \textup{(iv)}$).

\begin{enumerate}[label=\textup{(\roman*)},ref=\roman*]
\item\label{item:kp-def}
  Chirally Koszul (Definition~\ref{def:chiral-koszul-morphism}).
  \hfill \checkmark
\item\label{item:kp-pbw}
  PBW spectral sequence on $\barBgeom(\cA)$ collapses at~$E_2$
  (Theorem~\ref{thm:pbw-koszulness-criterion}).
  \hfill \checkmark
\item\label{item:kp-barcobar}
  Bar-cobar counit $\Omegach(\barBch(\cA)) \to \cA$ is a quasi-isomorphism
  (Theorem~\ref{thm:bar-cobar-inversion-qi}).
  \hfill \checkmark
\item\label{item:kp-hochschild}
  $\ChirHoch^*(\cA)$ polynomial, vanishing outside $\{0,1,2\}$
  (Theorem~\ref{thm:main-koszul-hoch}).
  \hfill \checkmark
\end{enumerate}

\medskip\noindent
\textbf{Tier~2: accessible conjectures.}

\begin{enumerate}[label=\textup{(\roman*)},ref=\roman*,resume]
\item\label{item:kp-ainfty}
  Minimal $A_\infty$-model has $m_n = 0$ for $n \ge 3$.
  Forward: Proposition~\ref{prop:ainfty-formality-implies-koszul};
  converse: Conjecture~\ref{conj:ainfty-koszul-characterization}.
  \hfill ?
\item\label{item:kp-ext}
  Ext diagonal: $\operatorname{Ext}^{p,q}_\cA(\omega_X,\omega_X)=0$ for $p \neq q$
  (Conjecture~\ref{conj:ext-diagonal-vanishing}).
  \hfill ?
\item\label{item:kp-barrbeck}
  Barr--Beck--Lurie comparison for $\barBch \dashv \Omegach$ is an equivalence.
  \hfill ?
\item\label{item:kp-facthom}
  Factorization homology $\textstyle\int_X \cA$ in degree~$0$ for all smooth~$X$.
  \hfill ?
\end{enumerate}

\medskip\noindent
\textbf{Tier~3: structural conjectures.}

\begin{enumerate}[label=\textup{(\roman*)},ref=\roman*,resume]
\item\label{item:kp-fmbdy}
  FM boundary acyclicity: $H^k(i_S^!\,\barBgeom_n(\cA)) = 0$ for $k \neq 0$, all strata~$S$.
  \hfill ?
\item\label{item:kp-nullvec}
  Kac--Shapovalov: $\det G_h \neq 0$ in the bar-relevant range
  (Conjecture~\ref{conj:kac-shapovalov-koszulness}).
  \hfill ?
\item\label{item:kp-dmod}
  $\mathcal{D}$-module purity: $\barBgeom_n(\cA)$ pure, $\operatorname{Ch}(\barBgeom_n) \subset \bigcup_S T^*_S\overline{C}_n(X)$.
  \hfill ?
\item\label{item:kp-lagrangian}
  Lagrangian complementarity: $\mathcal{M}_\cA, \mathcal{M}_{\cA^!} \hookrightarrow \mathcal{M}_{\mathrm{comp}}$ are $(-1)$-shifted Lagrangian (PTVV), conditional on perfectness.
  \hfill ?
\end{enumerate}


%% ------------------------------------------------------------------
\subsection{Operadic complexity}
\label{subsec:concordance-operadic-complexity}
\index{operadic complexity|textbf}
\index{shadow Postnikov tower!operadic complexity}

\begin{conjecture}[Operadic complexity;
  \ClaimStatusConjectured]
\label{conj:operadic-complexity}
\index{operadic complexity!conjecture}
For a chiral Koszul algebra~$\cA$, the three invariants
coincide: $r_{\max}(\cA) = d_\infty(\cA) = f_\infty(\cA)$,
where $r_{\max}$ is the shadow depth
(Theorem~\ref{thm:nms-finite-termination}), $d_\infty$ is
the $A_\infty$-depth (largest $n$ with $m_n^{\mathrm{tr}}
\neq 0$), and $f_\infty$ is the $L_\infty$-formality level
(largest $n$ with $\ell_n^{(0),\mathrm{tr}} \neq 0$ on
$\mathfrak{g}_\cA^{\mathrm{mod}}$).
\end{conjecture}

\noindent
The shadow--formality dictionary and archetype verification:

\begin{center}
\renewcommand{\arraystretch}{1.15}
\begin{tabular}{@{}llll@{}}
\toprule
\textbf{Shadow} & \textbf{$L_\infty$ bracket}
  & \textbf{Family} & $r_{\max}$ \\
\midrule
$\kappa$ (arity~$2$) & $\ell_2$
  & Heisenberg & $2$ \\
$\mathfrak{C}$ (arity~$3$) & $\ell_3^{\mathrm{tr}}$
  & Affine $\widehat{\fg}_k$ & $3$ \\
$\mathfrak{Q}$ (arity~$4$) & $\ell_4^{\mathrm{tr}}$
  & $\beta\gamma$ & $4$ \\
full tower & all $\ell_k^{\mathrm{tr}}$
  & $\mathrm{Vir}_c$ & $\infty$ \\
\bottomrule
\end{tabular}
\end{center}


%% ------------------------------------------------------------------
\subsection{Shadow depth classification}
\label{subsec:concordance-shadow-depth-classes}
\index{shadow depth!classification}
\index{Koszul property!shadow depth classes}

Shadow depth is orthogonal to Koszulness: all standard
families are chirally Koszul, yet their shadow depths differ.
See Theorem~\ref{thm:nms-archetype-trichotomy}.

\begin{center}
\renewcommand{\arraystretch}{1.15}
\begin{tabular}{@{}clll@{}}
\toprule
\textbf{Class} & \textbf{Name} & \textbf{Criterion}
  & \textbf{Examples} \\
\midrule
\textbf{G} & Gaussian
  & $r_{\max}=2$; formal
  & $\mathcal{H}_k$, free fermion, rank-$1$ lattice \\
\textbf{L} & Lie
  & $r_{\max}=3$; $\mathfrak{C}\neq 0$, $o_4=0$ (Jacobi)
  & $V_k(\fg)$, $Y(\fg)$ \\
\textbf{C} & Contact
  & $r_{\max}=4$; $\mathfrak{Q}\neq 0$, $o_5=0$ (abelian rigidity)
  & $\beta\gamma$, symplectic fermion\,? \\
\textbf{M} & Mixed
  & $r_{\max}=\infty$; non-formal
  & $\mathrm{Vir}_c$, $\mathcal{W}_N$ \\
\bottomrule
\end{tabular}
\end{center}

\noindent
The inclusions $\mathbf{G} \subset \mathbf{L} \subset
\mathbf{C} \subset \mathbf{M}$ are proper.  Shadow depth
classifies the complexity of~$\Theta_\cA$, not whether the
bar complex is acyclic
(Remark~\ref{rem:koszulness-vs-shadow-depth}).


%% ------------------------------------------------------------------
\subsection{Geometric, physics, and categorical
  characterizations}
\label{subsec:concordance-further-characterizations}
\index{Koszul property!geometric characterizations}
\index{Koszul property!physics characterizations}
\index{Koszul property!categorical characterizations}

\noindent\textbf{Geometric.}
\begin{itemize}[nosep]
\item \emph{Ran-space formality}: $E_2$-collapse $\Leftrightarrow$ genuswise bar-cobar formality $\Leftrightarrow$ $\tilde{m}_n^{(g)} = 0$ for $n \ge 3$ (Proposition~\ref{prop:e2-collapse-formality}); coherent extension to all $\operatorname{Ran}^{\le n}$ is open.
\item \emph{Tropical Koszulness}: acyclicity of $\barBgeom^{\mathrm{trop}}(\cA)$ for the planted-forest differential (Definition~\ref{def:planted-forest-coefficient-algebra}, Theorem~\ref{thm:planted-forest-tropicalization}).
\item \emph{Curve independence}: chiral Koszulness depends on formal-disk data, not on~$X$; higher-genus propagation governed by MK1--MK3 (Theorem~\ref{thm:genus-graded-koszul}) and the exceptional set~$\Sigma(\fg)$ (Theorem~\ref{thm:bar-cohomology-level-independence}).
\end{itemize}

\smallskip\noindent\textbf{Physics.}
\begin{itemize}[nosep]
\item \emph{One-loop exactness}: PBW $E_2$-collapse IS one-loop exactness (Remark~\ref{rem:koszulness-vs-shadow-depth}).
\item \emph{Kac--Shapovalov}: $\det G_h \neq 0$ in bar-relevant range $\Leftrightarrow$ Koszul (Conjecture~\ref{conj:kac-shapovalov-koszulness}); proved for $V_k(\fg)$ (Proposition~\ref{prop:pbw-universality}), open for $L_k(\fg)$ at admissible levels.
\item \emph{$C_2$-cofiniteness}: orthogonal to Koszulness; Koszul duality maps $C_2$-cofinite to non-$C_2$-cofinite.
\item \emph{Critical level}: $V_{-h^\vee}(\fg)$ is uncurved ($\kappa=0$), self-dual, center~$= \mathfrak{z}(\widehat{\fg})$ (Theorem~\ref{thm:critical-level-cohomology}).
\end{itemize}

\smallskip\noindent\textbf{Categorical.}
\begin{itemize}[nosep]
\item \emph{$t$-exactness}: $\barBch \colon \operatorname{FactMod}(\cA) \to \operatorname{FactCoMod}(\cA^!)$ exact for PBW/$\Omega$ $t$-structures (Proposition~\ref{prop:koszul-t-structures}).
\item \emph{Barr--Beck--Lurie}: effective descent for bar-cobar (Remark~\ref{rem:construction-vs-resolution}).
\item \emph{Semi-orthogonal decomposition}: weight-indexed SOD on $\operatorname{FactMod}(\cA)$ is split iff Koszul.
\item \emph{BGG}: bar resolution of $V_k(\fg)$-modules is linear; Koszul involution maps BGG to Cousin on $\operatorname{Gr}_G$ (Corollary~\ref{cor:bgg-koszul-involution}).
\item \emph{Convolution formality}: $\mathfrak{g}_\cA^{\mathrm{mod}}$ formal iff one-channel; stronger than Koszulness.
\item \emph{MC smoothness}: $\operatorname{MC}(\mathfrak{g}_\cA^{\mathrm{mod}})$ smooth at~$\Theta_\cA$ iff Koszul.
\end{itemize}


%% ------------------------------------------------------------------
\subsection{Bifunctor obstruction}
\label{subsec:concordance-bifunctor-obstruction}
\index{Robert-Nicoud--Wierstra!bifunctor obstruction}
\index{Koszul property!bifunctor obstruction}

The RNW no-bifunctor obstruction \cite{RNW19}
(\S\ref{sec:concordance-three-pillars}, identification
theorem~7) is structural: $\operatorname{hom}_\alpha(
\mathcal{C},\mathcal{A})$ extends to $\infty$-morphisms in
either slot separately but not both.  Koszulness does not kill
it; Koszulness makes it irrelevant by factoring all
constructions one slot at a time
(Theorem~\ref{thm:quillen-equivalence-chiral}).  This is why
the MC3 strategy
(Theorem~\ref{thm:mc3-type-a-resolution}) succeeds.
\label{rem:concordance-bifunctor-koszul}


%% ------------------------------------------------------------------
\subsection{Computational diagnostics}
\label{subsec:concordance-computational-diagnostics}
\index{Koszul property!computational diagnostics}
\index{bar generating function!algebraicity}

The bar generating function $\overline{P}_\cA(x,t)$ is
algebraic of degree~$2$ in~$t$ for all standard families with
$\dim\fg \le 8$, sharing the discriminant
$\Delta_\cA(x) = (1{-}3x)(1{+}x)$
across $\widehat{\mathfrak{sl}}_2$, $\beta\gamma$, and
Virasoro (Theorem~\ref{thm:ds-bar-gf-discriminant}).

\smallskip\noindent\textbf{Decision tree.}\enspace
(1) Pole order ($1 \Rightarrow$ G; $2 \Rightarrow$ L+).
(2) Freely generated $\Rightarrow$ Koszul (Proposition~\ref{prop:pbw-universality}).
(3) PBW + classical Koszulness of $\operatorname{gr}_F\cA$ (Theorem~\ref{thm:pbw-koszulness-criterion}).
(4) Bar Betti diagonal concentration (Corollary~\ref{cor:bar-cohomology-koszul-dual}).
(5) Generating function fit $\to$ discriminant.

\smallskip\noindent\textbf{Priority computations.}
\begin{itemize}[nosep]
\item $\mathcal{W}_4$ bar Betti numbers (Class~M + discriminant universality).
\item Bigraded Virasoro to $(n,h)=(4,16)$ (quartic resonance at GF level).
\item Symplectic fermion shadow (Class~C or~M assignment).
\end{itemize}


%% ------------------------------------------------------------------
\subsection{Cross-reference table}
\label{subsec:concordance-koszulness-xref}

\begin{center}
\renewcommand{\arraystretch}{1.15}
\begin{tabular}{@{}p{3.6cm}lp{4.5cm}@{}}
\toprule
\textbf{Item} & \textbf{Label} & \textbf{Location} \\
\midrule
PBW $E_2$-collapse
  & \texttt{thm:pbw-koszulness-criterion}
  & Ch.~\ref{chap:koszul-pairs} \\
Bar-cobar q.i.\
  & \texttt{thm:bar-cobar-inversion-qi}
  & Ch.~\ref{chap:bar-cobar-adjunction} \\
ChirHoch polynomial
  & \texttt{thm:main-koszul-hoch}
  & Ch.~\ref{chap:deformation-theory} \\
$A_\infty$ ($\Rightarrow$)
  & \texttt{prop:ainfty-formality-implies-koszul}
  & Ch.~\ref{chap:koszul-pairs} \\
$A_\infty$ ($\Leftarrow$)
  & \texttt{conj:ainfty-koszul-characterization}
  & Ch.~\ref{chap:koszul-pairs} \\
Ext diagonal
  & \texttt{conj:ext-diagonal-vanishing}
  & Ch.~\ref{chap:koszul-pairs} \\
Kac--Shapovalov
  & \texttt{conj:kac-shapovalov-koszulness}
  & Ch.~\ref{chap:koszul-pairs} \\
Deformation quant.
  & \texttt{rem:dq-koszul-compatibility}
  & Ch.~\ref{chap:koszul-pairs} \\
Meta-theorem
  & \texttt{thm:koszul-equivalences-meta}
  & Ch.~\ref{chap:koszul-pairs} \\
Operadic complexity
  & \texttt{conj:operadic-complexity}
  & \S\ref{subsec:concordance-operadic-complexity} \\
Archetype classification
  & \texttt{thm:nms-archetype-trichotomy}
  & App.~\ref{app:nonlinear-modular-shadows} \\
Finite termination
  & \texttt{thm:nms-finite-termination}
  & App.~\ref{app:nonlinear-modular-shadows} \\
Discriminant
  & \texttt{thm:ds-bar-gf-discriminant}
  & Landscape census \\
Envelope-shadow
  & \texttt{def:envelope-shadow-functor}
  & \S\ref{sec:concordance-nishinaka-vicedo} \\
Finite-jet rigidity
  & \texttt{prop:finite-jet-rigidity}
  & \S\ref{sec:concordance-nishinaka-vicedo} \\
Gaussian collapse
  & \texttt{prop:gaussian-collapse-abelian}
  & \S\ref{sec:concordance-nishinaka-vicedo} \\
\bottomrule
\end{tabular}
\end{center}

\noindent
\emph{Status}: four equivalences proved
(\ref{item:kp-def})$\Leftrightarrow$(\ref{item:kp-pbw})$\Leftrightarrow$(\ref{item:kp-barcobar})$\Leftrightarrow$(\ref{item:kp-hochschild});
eight conjectural.  Most accessible next target:
(\ref{item:kp-ainfty}), requiring the HPL contracting
homotopy on the formal disk.  Most ambitious:
(\ref{item:kp-lagrangian}), identifying Koszulness with
$(-1)$-shifted Lagrangian transversality.


%% ====================================================================
%%  NON-PRINCIPAL W-ALGEBRA DUALITY: THE TRANSPORT FRONTIER
%% ====================================================================
\section{Non-principal \texorpdfstring{$\mathcal{W}$}{W}-algebra
duality: the transport frontier}
\label{sec:concordance-non-principal-w}
\index{W-algebra@$\mathcal{W}$-algebra!non-principal duality!concordance}

After the principal Feigin--Frenkel case
(Theorem~\ref{thm:w-algebra-koszul-main}), the non-principal
duality problem splits into two axes
(Principle~\ref{princ:transport-orbit-factorization}):
\emph{transport along the reduction network} and
\emph{identification of the dual orbit/level}.

\paragraph{Proved core.}
\begin{itemize}
\item \emph{Principal}: Feigin--Frenkel duality, all simple~$\fg$.
\item \emph{Subregular/minimal in $\mathfrak{sl}_3$}:
  Bershadsky--Polyakov duality (Proposition~\ref{prop:bp-duality}).
\item \emph{Hook-type corridor in type~$A$}:
  full hook-type duality at generic level
  (Theorem~\ref{thm:hook-transport-corridor}), with inverse
  reductions connecting all hook-type nodes
  \cite{FehilyHook,CLNS24,CFLN24,GenraStages,ButsonNair}.
\item \emph{Transport propagation}: duality propagates from hook
  vertices through the reduction graph~$\Gamma_N$ to every
  partition in the transport-closure
  (Proposition~\ref{prop:transport-propagation}).
\end{itemize}

\paragraph{Active frontier.}
\begin{itemize}
\item \emph{Type-$A$ transport-to-transpose}
  (Conjecture~\ref{conj:type-a-transport-to-transpose}):
  the transport-closure of the hook vertices equals
  $\mathrm{Par}(N)$ at generic level.
\item \emph{General orbit duality}
  (Conjecture~\ref{conj:w-orbit-duality}):
  extends to all simple~$\fg$ via Barbasch--Vogan/Spaltenstein
  duality.
\item \emph{DS--KD intertwining for arbitrary nilpotent}
  (Conjecture~\ref{conj:ds-kd-arbitrary-nilpotent}):
  bar--cobar duality commutes with DS reduction at arbitrary~$f$.
\end{itemize}

\paragraph{Key constraint.}
The difficulty is \emph{not} the existence of quantum DS
reduction for arbitrary~$f$ (Kac--Roan--Wakimoto~\cite{KRW});
it is proving that bar--cobar/Koszul duality commutes with
arbitrary-nilpotent reduction, identifying the correct dual
nilpotent, and determining the non-principal level transform.


\section{Bridges to Volume~II}
\label{sec:cross-volume-bridges}
\index{Volume II!bridges from Volume I}

Volume~II develops the $A_\infty$-chiral Hochschild theory in
three space-time dimensions with the operad
$\mathrm{SC}^{\mathrm{ch,top}}$.  Homotopy--Koszulity of
$\mathrm{SC}^{\mathrm{ch,top}}$ is proved in Volume~II
(via Kontsevich formality and transfer from the classical
Swiss--cheese operad), so the bar--cobar adjunction for
$\mathrm{SC}^{\mathrm{ch,top}}$-algebras is a Quillen
\emph{equivalence}.  Volume~II has \emph{no remaining conjectural
algebraic inputs}; its only conditional inputs are the
standing analytic hypotheses (H1)--(H4).
Five bridges connect the two volumes:

\begin{proposition}[Bar-cobar bridge; \ClaimStatusProvedElsewhere]
\label{conj:vol2-bar-cobar-bridge}
The $\mathrm{SC}^{\mathrm{ch,top}}$ bar-cobar Quillen equivalence of
Volume~II recovers the chiral bar-cobar adjunction of Theorem~A
(Theorem~\ref{thm:bar-cobar-isomorphism-main}) for boundary
operators of the 3d theory.  The closed color specializes to the BD
chiral operad (Vol~II, associated graded proposition).
\end{proposition}

\begin{conjecture}[Hochschild bridge; \ClaimStatusConjectured]
\label{conj:vol2-hochschild-bridge}
The bulk algebra of the 3d theory is the chiral Hochschild cochain
complex $C^*_{\mathrm{ch}}(\cA)$, recovering Theorem~H
(Theorem~\ref{thm:w-algebra-hochschild}) via dimensional
descent.
\end{conjecture}

\begin{proposition}[DK-YBE bridge; \ClaimStatusProvedElsewhere]
\label{conj:vol2-dk-ybe-bridge}
The spectral braiding $r(z) = \int_0^\infty e^{-pz}
\{a_\lambda b\}\,dp$ (the Laplace transform of the
$\lambda$-bracket) recovers the DK-0 chain-level
monodromy comparison (Theorem~\ref{thm:derived-dk-affine}).
\end{proposition}

\begin{proposition}[\texorpdfstring{$\mathcal{W}$}{W}-algebra bridge at genus $0$; \ClaimStatusProvedElsewhere]
\label{conj:vol2-w-algebra-bridge}
The Feynman $m_k$ operations of the 3d theory, restricted to
the $\mathcal{W}$-algebra family, match the bar differential computed in the
$\mathcal{W}$-algebra chapters of Part~II at genus~$0$.
Extension to genus~$1$ is proved
(Volume~II, Theorem~\ref{thm:mc5-genus-one-bridge}:
$d_{\mathrm{fib}}^2 = \kappa(\cA)\cdot\omega_1$ via the Arnold
defect on the torus); extension to $g \geq 2$ remains conjectural
(MC5, requiring Costello renormalization).
\end{proposition}

\begin{conjecture}[BV functor bridge; \ClaimStatusConjectured]
\label{conj:vol2-bv-functor-bridge}
Analytic hypotheses (H1)--(H4) of Volume~II define the
physics-to-algebra functor that realizes the
$\mathrm{BV} = \mathrm{bar}$ identification at higher
genus (MC5, Conjecture~\ref{conj:master-bv-brst}).
\end{conjecture}


\section[Towards modular homotopy theory]{Towards modular homotopy theory for factorization algebras on curves}
\label{sec:modular-koszul-programme}
\index{modular Koszul duality!project}

The modular characteristic theorem (Theorem~D,
Theorem~\ref{thm:modular-characteristic}) establishes the scalar
logarithmic datum~$\kappa(\cA)$.  This section identifies the six
components---five structural extensions and the Hochschild
cohomology theorem---needed to pass from the proved modular Koszul
core to modular homotopy theory for factorization algebras on
curves---that is, to complete the nilpotence-periodicity
correspondence
(Remark~\ref{rem:nilpotence-periodicity}) from the scalar level to
the full categorical level:

\begin{center}
\renewcommand{\arraystretch}{1.3}
\begin{tabular}{@{}cp{4.4cm}p{3.8cm}p{4.6cm}@{}}
\toprule
\textbf{Label} & \textbf{Statement} & \textbf{Status} & \textbf{Reference} \\
\midrule
$A_{\mathrm{mod}}$ & Bar-cobar intertwined with Verdier, & \textbf{Proved} & Thm~\ref{thm:bar-cobar-isomorphism-main} \\
  & functorial over $\overline{\mathcal{M}}_{g,n}$ & & \\[2pt]
$B_{\mathrm{mod}}$ & Inversion on Koszul locus; coderived & Proved on Koszul locus; & Thm~\ref{thm:higher-genus-inversion}, \\
  & persistence off it & coderived persistence conjectural & \S\ref{subsec:coderived-ran} \\[2pt]
$C_{\mathrm{mod}}$ & Shifted symplectic complementarity & \textbf{Proved} & Thms~\ref{thm:quantum-complementarity-main}, \\
  & (PTVV Lagrangian embedding) & & \ref{thm:shifted-symplectic-complementarity} \\[2pt]
Index & GRR: genus series $= \kappa(\cA) \cdot (\hat{A}(ix) - 1)$ & \textbf{Proved} & Thm~\ref{thm:family-index} \\[2pt]
DK & DK-0/1/$1\frac{1}{2}$: chain-level, eval-locus, lattice; & DK-0/1/$1\frac{1}{2}$: proved (all types); & Thms~\ref{thm:derived-dk-affine}, \\
  & DK-2/3: factorization DK; & DK-2/3: \textbf{proved} (eval-gen.\ core, all types, domain: $\mathcal{O}_{\mathrm{poly}}$; & \ref{thm:derived-dk-yangian}, \ref{thm:factorization-dk-eval}, \\
  & DK-4/5: dg-shifted/triple bridge & uses Molev PBW \cite{molev-yangians}); DK-4: ML proved, alg.\ id.\ open; DK-5: conj. & \ref{thm:dk-fd-typeA}, \ref{cor:dk23-all-types}, \ref{thm:rtt-mittag-leffler}; Conj.~\ref{conj:full-dk-bridge} \\[2pt]
$H$ & Chiral Hochschild: $\ChirHoch^*(\Walg^k(\fg)) \cong \mathbb{C}[\Theta_1,\ldots,\Theta_r]$; & \textbf{Proved} (generic level, & Thms~\ref{thm:w-algebra-hochschild}, \\
  & period $2$ (rank $1$), polynomial (rank $\geq 2$); & critical level: ProvedElsewhere) & \ref{thm:virasoro-hochschild}, \ref{thm:critical-level-cohomology}; \\
  & Koszul functoriality via Connes~$S$ & & Thm~\ref{thm:main-koszul-hoch} \\
\bottomrule
\end{tabular}
\end{center}

\begin{remark}[Four levels of modular invariant]%
\label{rem:four-levels}%
\index{modular homotopy theory!four levels}%
The six components above implement a hierarchy of modular
invariants of increasing categorical depth.
The proved scalar level is the lowest;
the full modular homotopy type is the highest:
\begin{center}
\renewcommand{\arraystretch}{1.3}
\begin{tabular}{@{}clll@{}}
\toprule
\textbf{Level} & \textbf{Invariant} & \textbf{Type} & \textbf{Status} \\
\midrule
$0$ & $\kappa(\cA) \cdot \hat{A}$-genus
  & Numbers $F_g$ & \textbf{Proved} (Thm~D) \\[2pt]
$1$ & $H^*(\barB^{(g)}(\cA), \Dg{g})$
  & Graded vector spaces & \textbf{Resolved} at integrable levels \\
  & & & (Thm~\ref{thm:bar-cohomology-level-independence}) \\[2pt]
$2$ & Variation of $H^*$ over $\mathcal{M}_g$
  & Flat connection (KZ/Hitchin) & Open \\[2pt]
$3$ & Factorization and sewing
  & Chain-level modular functor & Structural \\
  & & (Thm~\ref{thm:chain-modular-functor}) & \\
\bottomrule
\end{tabular}
\end{center}
Theorem~D is the \emph{integrability condition}:
$\mathrm{obs}_g = \kappa \cdot \lambda_g$ is a pure
tautological class, free of boundary or non-tautological
contributions, guaranteeing that the curvature of the genus-$g$
bar complex is absorbable by a period correction.
The corrected complex $(\barB^{(g)}, \Dg{g})$ with
$\Dg{g}^{\,2} = 0$ (Level~$1$) is the
\emph{derived fiber} of the factorization homology bundle
$\int_{\Sigma_g}\!\cA$ over~$\mathcal{M}_g$; its variation
(Level~$2$) recovers the monodromy representation of the mapping
class group; its factorization structure (Level~$3$) recovers the
full modular functor
(Theorem~\ref{thm:chain-modular-functor}).  The $\hat{A}$-genus
is the \emph{first Chern class} of this bundle---its Euler
characteristic, not the bundle itself.

\smallskip\noindent
\emph{Smoking gun}: for $\widehat{\mathfrak{sl}}_2$ at level~$k$,
$\kappa = 3(k{+}2)/4$ is one number, yet the Verlinde formula
gives a bundle of rank $k{+}1$ at genus~$1$
(equation~\eqref{eq:verlinde-general}), growing with~$g$.
Bar cohomology $H^0(\barB^{(1)}(\widehat{\mathfrak{sl}}_{2,k}),
\Dg{1})$ should recover the Verlinde dimension---a concrete,
computable test that Level~$1$ carries representation-theoretic
data invisible to Level~$0$.  See
\S\ref{subsec:chain-modular-functor}.

\smallskip\noindent
\emph{Resolution of genus-independence at integrable levels.}
Theorem~\ref{thm:bar-cohomology-level-independence} proves that
bar cohomology dimensions are polynomial in $\lambda = k + h^\vee$,
constant outside a finite exceptional set~$\Sigma_n$.  For
$\mathfrak{sl}_2$ at degree~$2$: $\Sigma_2 = \{0\}$ (critical level
only).  Since integrable levels $k \ge 1$ give $\lambda \ge 3 > 0$,
they are \emph{generic}---not exceptional---and the fiberwise bar
cohomology is genus-independent there
(Proposition~\ref{prop:integrable-level-independence}).  The
integer conformal block dimension $\dim V_{g,k}$ relates to the
normalized partition function~$Z_g$ via $\dim V = Z_g \cdot
S_{00}^{\,2-2g}$ (Remark~\ref{rem:two-verlinde-normalizations}).
\end{remark}

\noindent
The subsections below develop each component, separating what is
established from what remains.  For $A_{\mathrm{mod}}$: the present
proof establishes functorial compatibility of bar and cobar over
$\overline{\mathcal{M}}_{g,n}$; whether this lifts to an equivalence
of modular operads is a further question, not addressed here.

\medskip

\subsection{The coderived Ran-space formalism}
\label{subsec:coderived-ran}

The bar-cobar adjunction is most naturally an equivalence of
$\infty$-categories---the categorical $\exp/\log$ correspondence
(Remark~\ref{rem:nilpotence-periodicity}):
\[
\bar{B}_X \colon
\operatorname{Alg}_{\mathrm{aug}}(\operatorname{Fact}(X))
\;\rightleftarrows\;
\operatorname{CoAlg}_{\mathrm{conil}}(\operatorname{Fact}(X))
\,:\!\Omega_X,
\]
where $\operatorname{Fact}(X)$ denotes factorization algebras on
$\operatorname{Ran}(X)$ (cf.\ Appendix~\ref{app:existence-criteria}).

On the Koszul locus this is an equivalence: $\exp \circ \log = \mathrm{id}$
on the convergence domain.  When curvature is present---that is,
when the categorical logarithm acquires non-trivial monodromy---the
correct ambient category is not the ordinary derived category but the
\emph{coderived category} in the sense of
Positselski~\cite{Positselski11}: the analytic continuation that
extends the $\exp/\log$ correspondence past the convergence radius.

Curved objects produce acyclic-but-not-contractible complexes that the
derived category erases too brutally; the coderived category retains
the monodromy data that the ordinary derived category discards, just
as the multi-valued logarithm retains branch information that a
single-valued determination loses.

The provisional coderived category
$\mathsf{D}^{\mathrm{co}}_{\mathrm{prov}}(X)$
(Definition~\ref{def:provisional-coderived},
Proposition~\ref{prop:coderived-adequacy})
provides a working substitute within the scope of this monograph.
The Koszul locus---where bar-cobar inversion
holds unconditionally (Theorem~\ref{thm:higher-genus-inversion})---is
the locus where the coderived and derived categories agree.

\emph{What is established.}
Appendix~\ref{app:existence-criteria} compares the strict,
completed, and literature existence paradigms for bar-cobar on
$\operatorname{Ran}(X)$.  Theorem~\ref{thm:bar-cobar-isomorphism-main}
provides the adjunction; Theorem~\ref{thm:higher-genus-inversion}
provides inversion on the Koszul locus.
Appendix~\ref{app:coderived},
\S\ref{subsec:coderived-ran-formalism} develops the coderived
and contraderived categories on~$\operatorname{Ran}(X)$: curved
factorization coalgebras with CDG-factorization compatibility
(Definition~\ref{def:curved-fact-coalgebra}), coacyclic
factorization objects
(Definition~\ref{def:coacyclic-fact}), stratified conservative
restriction via the injective characterization
(Theorem~\ref{thm:stratified-conservative-restriction}),
and the factorization co-contra correspondence
(Theorem~\ref{thm:fact-co-contra-general}).  The provisional
coderived category embeds fully faithfully
(Proposition~\ref{prop:provisional-embedding}).

\emph{Now complete.}
The formerly conjectural off-Koszul bar-cobar inversion on
$\operatorname{Ran}(X)$ is proved
(Theorem~\ref{thm:off-koszul-ran-inversion}): the counit
$\Omega_X \barB_X(\cA) \to \cA$ is an isomorphism in
$D^{\mathrm{co}}(\mathrm{CoFact}(X))$ without Koszulness.
The proof follows the same three-step assembly as
Theorem~\ref{thm:fact-co-contra-general}: (i) on each stratum
$\operatorname{Ran}_n$, Positselski~\cite[Theorem~7.2]{Positselski11}
gives the coderived equivalence for the restricted CDG-coalgebra;
(ii)~factorization compatibility of the bar-cobar counit;
(iii)~conservative assembly via
Theorem~\ref{thm:stratified-conservative-restriction}.
This closes structural gap~H1 and completes Future~2.

\subsection{The Lagrangian form of complementarity}

Theorem~C (Theorem~\ref{thm:quantum-complementarity-main}) establishes
the direct-sum decomposition
$Q_g(\cA) \oplus Q_g(\cA^!)
= H^*(\overline{\mathcal{M}}_g, \mathcal{Z}(\cA))$
via the eigenspace decomposition of the Verdier involution.
The natural strengthening is a Lagrangian-polarization theorem,
now proved:

\begin{theorem}[Lagrangian complementarity;
\ClaimStatusProvedHere]\label{thm:lagrangian-complementarity}
\index{Lagrangian!complementarity|textbf}
For a modular Koszul pair $(\cA, \cA^!)$,
the complex
$C_g := R\Gamma(\overline{\mathcal{M}}_g, \mathcal{Z}_{\cA})$
carries two shifted symplectic structures, each admitting
$Q_g(\cA)$ and $Q_g(\cA^!)$ as complementary Lagrangians in the
sense of Pantev--To\"en--Vaqui\'e--Vezzosi~\cite{PTVV13}:
\begin{enumerate}[label=\textup{(\roman*)}]
\item a $(-(3g{-}3))$-shifted symplectic structure from
  Verdier duality on the center local system
  \textup{(}Proposition~\textup{\ref{prop:ptvv-lagrangian}}\textup{)};
\item a $(-1)$-shifted symplectic structure on the formal moduli
  $\mathrm{Def}_g(\cA)$ from the BV antibracket
  \textup{(}Theorem~\textup{\ref{thm:shifted-symplectic-complementarity}}\textup{)}.
\end{enumerate}
\end{theorem}

\begin{remark}[Degree clarification]\label{rem:lagrangian-conj-scope}
The original formulation of this result stated a single
$(-1)$-shifted symplectic structure on~$C_g$.  The proved version
reveals two structures of different degrees: the Verdier pairing
gives $(-(3g{-}3))$-shifted on~$C_g$ itself, while the BV
antibracket gives $(-1)$-shifted on the formal moduli.  Both support
the same Lagrangian polarization.
See Remark~\ref{rem:ptvv-relation} for the compatibility.
\end{remark}

\subsection{The analytic completion programme}
\label{subsec:analytic-completion-programme}
\index{analytic completion!programme}

The proved algebraic engine operates in the world of formal
factorization algebras, cochain complexes, and dg categories.  A
genuine local-to-global field theory also requires analytic sewing
(convergent partition functions, Hilbert-space completions, and
positivity).  Recent work of
Moriwaki~\cite{Moriwaki26a, Moriwaki26b} constructs conformally flat
factorization homology valued in $\operatorname{IndHilb}$ by left Kan
extension from conformally flat disk algebras; for the Heisenberg
VOA, the conformally flat 2-disk algebra on
$\operatorname{Sym} A^2(D)$ (the symmetric algebra of the Bergman
space) is identified with the ind-Hilbert completion of the affine
Heisenberg vertex algebra.  Adamo--Moriwaki--Tanimoto~\cite{AMT24}
prove conformal Osterwalder--Schrader axioms for a class of unitary
full VOAs under energy/spectral growth and local $C_1$-cofiniteness.

The programme aims to bridge the algebraic engine to these analytic
structures.  Its objects are the \emph{sewing envelope}
$\cA^{\mathrm{sew}}$ (the Hausdorff completion for the locally convex
topology generated by sewing-amplitude seminorms), the \emph{analytic
bar coalgebra} $\barB^{\mathrm{an}}(\cA)$ (graph-norm closure of the
algebraic bar), and the \emph{analytic shadow}
$Q_g^{\mathrm{an}}(\cA) := D^{\mathrm{co}}(
\barB_g^{\mathrm{an}}(\cA)\text{-comod})$.  At genus~$g$ with
curvature $\mcurv{g} \neq 0$, the coderived passage of
\S\ref{subsec:coderived-ran} is mandatory: ordinary cohomology fails,
and the correct receptacle is the coderived analytic category.

Four target theorems define the programme
(Chapter~\ref{ch:genus-complete},
\S\ref{sec:analytic-completion}):

\begin{enumerate}[label=\textup{(\Alph*$_{\mathrm{an}}$)}]
\item \emph{Heisenberg sewing theorem}
  (Theorem~\ref{thm:heisenberg-sewing}): the sewing envelope
  of the algebraic Heisenberg VOA is exactly
  $\operatorname{Sym} A^2(D)$; the completed bar differential is
  the closure of the Gaussian collision operator on Bergman tensors;
  every closed amplitude is a Fredholm determinant on the
  one-particle Bergman space.
\item \emph{Lattice sewing envelope}
  (Conjecture~\ref{conj:lattice-sewing}): every positive-definite
  even lattice equipped with an analytic theta-datum admits a
  canonical sewing envelope with charge-refined Hilbert sectors.
\item \emph{Analytic realization criterion}
  (Conjecture~\ref{conj:analytic-realization}): a unitary full VOA
  satisfying conformal OS, polynomial spectral growth, and
  HS-sewing admits a sewing envelope, a conformally flat 2-disk
  algebra, and a higher-genus coderived shadow.
\item \emph{Boundary bar duality}
  (Conjecture~\ref{conj:boundary-bar-duality}): for an
  analytically Koszul~$\cA$, completed boundary modules at
  genus~$0$ are equivalent to analytic comodules over
  $\barB^{\mathrm{an}}(\cA)$; at curved genus, the correct
  objects are analytic contramodules over the curved dual.
\end{enumerate}

\noindent
\emph{Status.}  HS-sewing is proved for the standard landscape
(Theorem~\ref{thm:general-hs-sewing}); the Heisenberg sewing
theorem is proved
(Theorem~\ref{thm:heisenberg-sewing}).
The sewing envelope universal property, analytic conilpotency,
and boundary bar duality are conjectural;
these concern the categorical equivalence of analytic module
categories beyond convergence of partition functions.

\subsection{The universal Maurer--Cartan class}

The genus tower $\{\barB^{(g)}(\cA)\}_{g \geq 0}$
(Theorem~\ref{thm:master-tower}) is controlled by a universal
obstruction.

\begin{theorem}[Universal MC class;
\ClaimStatusProvedHere]\label{thm:universal-MC}
\index{Maurer--Cartan class!universal}
For a modular Koszul chiral algebra $\cA$ with simple Lie
symmetry~$\mathfrak{g}$, there exists a class
\[
\Theta_{\cA} \;\in\;
\operatorname{MC}\!\left(
  \operatorname{Def}_{\mathrm{cyc}}(\cA)
  \;\widehat{\otimes}\;
  R\Gamma\!\left(\overline{\mathcal{M}}_{g,\bullet},\, \mathbb{Q}\right)
\right)
\]
controlling the full genus tower as a single Maurer--Cartan deformation,
with three properties:
\begin{enumerate}[label=\textup{(\roman*)}]
\item Its image under the scalar trace is the sequence of obstruction
  coefficients: $\operatorname{tr}(\Theta_{\cA})
  = \sum_{g \geq 1} \kappa(\cA) \cdot \lambda_g$.
\item It is compatible with clutching:
  $\xi^*(\Theta_{\cA})
  = \Theta_{\cA} \star \Theta_{\cA}$
  under the boundary maps
  $\xi \colon \overline{\mathcal{M}}_{g_1,n_1+1} \times
  \overline{\mathcal{M}}_{g_2,n_2+1} \to
  \overline{\mathcal{M}}_{g,n}$.
\item Under Verdier duality:
  $\mathbb{D}(\Theta_{\cA}) = \Theta_{\cA^!}$.
\end{enumerate}
\end{theorem}

\begin{proof}
This is Theorem~\ref{thm:universal-theta}, proved via the
assembly of Theorem~\ref{thm:mc2-full-resolution}: the cyclic
$L_\infty$-algebra $\Defcyc(\cA)$ is constructed by
Theorem~\ref{thm:cyclic-linf-graph}, the completed tensor product
converges with clutching by
Proposition~\ref{prop:geometric-modular-operadic-mc}, and the
tautological line support $o_g = \kappa \lambda_g$ is established
by Theorem~\ref{thm:tautological-line-support}.
\end{proof}

\begin{remark}[Scope]\label{rem:universal-MC-scope}
The scalar shadow $\kappa(\cA) \cdot \lambda_g$ is established
(Theorem~\ref{thm:genus-universality}).  The clutching compatibility
is implicit in the modular operad structure of the bar complex
(Theorem~\ref{thm:prism-higher-genus}).  At the one-channel level,
Proposition~\ref{prop:tautological-line-support-criterion} shows that
these two geometric inputs would already force tautological-line
support once they separate the relevant sector, and
Proposition~\ref{prop:one-channel-verdier-criterion} reduces the
remaining duality flank to a $\sigma$-stable Lagrangian two-plane on
that sector.  Proposition~\ref{prop:one-channel-ptvv-criterion} then
reduces the remaining H-level lift to perfect one-channel projector
subcomplexes whose cross Verdier pairing is a PTVV quasi-isomorphism,
and Proposition~\ref{prop:one-channel-chain-model-criterion} reduces
that lift further to explicit one-channel coderivation subcomplexes in
$\Defcyc(\cA)$ and $\Defcyc(\cA^!)$.  Proposition~\ref{prop:one-channel-seed-criterion}
then reduces those subcomplexes further to finite low-bar-length
cyclic coderivation seeds in
$\operatorname{CoDer}^{\mathrm{cyc}}(\widehat{\barB}_X(\cA))[1]$ and
$\operatorname{CoDer}^{\mathrm{cyc}}(\widehat{\barB}_X(\cA^!))[1]$.
Proposition~\ref{prop:one-channel-minimal-seed-packet-criterion}
compresses those finite seed sets further to one distinguished
degree-$2$ cocycle on each side together with bar-length-$\leq 3$
correction packets and a finite pairing matrix.
Proposition~\ref{prop:one-channel-visible-lowarity-packet-criterion}
then identifies the remaining packet with the three visible low-arity
ingredients already supported by the current MC2 surface: the
simple-pole bracket sector, the normalized double-pole pairing matrix,
and the Killing $l_3^\eta$ sector.
Proposition~\ref{prop:one-channel-canonical-transfer-criterion} then
reduces that visible packet further to a canonical transfer package on
the one-channel seed spaces: one cyclic seed, one shared
generator-seed lift producing the Killing $l_3^\eta$ sector, and one
functorial normalization splitting off the one-channel cocycle line.
Proposition~\ref{prop:one-channel-transfer-law-criterion} then reduces
that canonical package further to an explicit root-string transfer law
on those seed spaces: the genus-$2$/genus-$3$ profile
\[
O_2=t^2(ae_{12}+bf_1-mf_2),
\qquad
O_3^\eta=a m t^3\eta,
\]
together with obstruction-side recovery of $(a,b,m)$ and the same
coefficient-extraction normalization.
Proposition~\ref{prop:one-channel-root-string-chart-criterion} then
reduces that explicit law further to a chart-existence/uniqueness
statement: the ordered root-string chart is forced by the one-channel
support graph and normalization data, up to rescaling of the three seed
lines.
Proposition~\ref{prop:one-channel-intrinsic-line-detection-criterion}
then reduces that chart statement further to intrinsic line detection:
the seed, simple-pole, and Killing support lines must be canonically
picked out by $l_{2,\mathrm{sp}}$, $\beta_g^{\mathrm{dp}}$,
$\nu_g^{\mathrm{sp}}$, and the support of
$(O_{2,g},O_{3,g}^\eta)$.
Proposition~\ref{prop:one-channel-automorphism-rigidity-criterion}
then reduces that remaining line-detection step to automorphism
rigidity: those support lines are the unique invariant lines for the
one-channel support automorphism group compatible with the bracket,
pairing, normalization, and obstruction profile.
Proposition~\ref{prop:one-channel-support-graph-stabilizer-criterion}
then reduces that rigidity statement further to a finite stabilizer
computation on the labeled one-channel support graph: the fixed
vertices must be exactly the simple-pole, residual-kernel, and Killing
support lines.
Proposition~\ref{prop:one-channel-incidence-orbit-criterion}
then reduces that stabilizer statement further to an explicit
incidence-matrix / orbit-count computation on the visible one-channel
graph: the relevant support lines are exactly the singleton orbits of
the permutation group determined by the bracket, pairing,
normalization, and support indicators.
Proposition~\ref{prop:one-channel-visible-orbit-table-criterion}
then reduces that bounded permutation problem further to a universal
three-case visible orbit table on the
$m=1,2,3$ root-string types, equivalently
$\mathfrak{sl}_3,\mathfrak{sp}_4,\mathfrak{g}_2$.
Proposition~\ref{prop:one-channel-canonical-universal-orbit-table-criterion}
then collapses that three-case verification to direct
lookup/identification against one fixed labeled orbit table.
Proposition~\ref{prop:one-channel-universal-invariant-signature-criterion}
then collapses that fixed-table identification further to the minimal
visible invariant signature packet forcing that table.
Proposition~\ref{prop:one-channel-signed-seed-character-criterion}
then collapses that invariant packet further to the universal signed
seed-character law recovering it.
Proposition~\ref{prop:one-channel-two-sign-plus-normalization-scalar-criterion}
then collapses that character further to the reduced datum of two
root-string signs and one normalization scalar.
Proposition~\ref{prop:one-channel-parity-sign-plus-normalization-scalar-criterion}
then collapses that reduced datum further to one root-string parity
sign plus one normalization scalar.
Proposition~\ref{prop:one-channel-parity-forcing-criterion}
then collapses that parity-sign datum further to one chart-normalized
seed scalar.
The scalar-saturated class $\Theta_{\cA}$ is now constructed
for one-channel families
(Theorem~\ref{thm:mc2-full-resolution}): (a)~the cyclic
deformation complex $\operatorname{Def}_{\mathrm{cyc}}(\cA)$
exists as a cyclic $L_\infty$ algebra
(Theorem~\ref{thm:cyclic-linf-graph}); (b)~the Maurer--Cartan
equation in the completed tensor product with tautological classes
is solved at the scalar level by the conjunction of
Proposition~\ref{prop:geometric-modular-operadic-mc} and
Theorem~\ref{thm:tautological-line-support}.
The non-scalar content is developed through the shadow
Postnikov tower
(Definition~\ref{def:shadow-postnikov-tower}): finite-order
shadows through arity~$4$ are proved constructively, and the
all-arity convergence is guaranteed by the bar-intrinsic
construction (Theorem~\ref{thm:mc2-bar-intrinsic};
Theorem~\ref{thm:recursive-existence}).
Scalar saturation (one-channel universality) is proved
for all algebras in algebraic families with rational OPE
coefficients (Theorem~\ref{thm:algebraic-family-rigidity}),
covering the entire standard Lie-theoretic landscape at
all non-critical levels including admissible ones.
The proof bypasses module-category semisimplicity via
Whitehead reduction to a finite-dimensional linear system
and algebraic semicontinuity of rank.

\emph{Homotopy template} (Convention~\ref{conv:hms-levels}): Type~I
(existence of a homotopy object).  At H-level: a filtered cyclic $L_\infty$
algebra $\operatorname{Def}_{\mathrm{cyc}}(\cA)$ and an MC element $\Theta_\cA$
with trace, clutching, and Verdier compatibility.  The scalar shadow is
$\kappa(\cA) \cdot \lambda_g$ (proved, Theorem~\ref{thm:genus-universality}).
The full non-scalar class is Theorem~\ref{thm:universal-theta}
(Theorem~\ref{thm:master-theta}).
\end{remark}

\begin{conjecture}[Categorical modular Koszul duality;
\ClaimStatusConjectured]%
\label{conj:categorical-modular-kd}%
\index{modular Koszul duality!categorical}%
For each genus~$g$ and each modular Koszul chiral
algebra~$\cA$, the curved bar-cobar adjunction at genus~$g$
induces an equivalence
\[
\mathsf{D}^{\mathrm{co}}
\bigl(\barB^{(g)}(\cA)\text{-}\mathsf{comod}\bigr)
\;\simeq\;
\mathsf{D}^b\bigl(\cA\text{-}\mathsf{mod}^{\Sigma_g}\bigr)
\]
between the \emph{coderived category}
\textup{(}\S\textup{\ref{subsec:coderived-ran})} of curved
$\barB^{(g)}(\cA)$-comodules and the bounded derived category of
$\cA$-modules twisted by the Hodge data of~$\Sigma_g$.
This equivalence varies over $\mathcal{M}_g$ as a flat functor:
\begin{enumerate}[label=\textup{(\roman*)}]
\item At $g = 0$, it reduces to bar-cobar inversion
  \textup{(}Theorem~\textup{\ref{thm:higher-genus-inversion})}.
\item Applied to integrable representations of
  $\widehat{\fg}_k$, the functor $H^0$ recovers the
  Tsuchiya--Ueno--Yamada space of conformal blocks
  \textup{(}Remark~\textup{\ref{rem:chain-vs-classical-mf})},
  and its monodromy gives the KZ/Hitchin connection.
\item The scalar shadow
  $c_1(\det \mathcal{V}_g) = \kappa(\cA) \cdot \lambda_g$
  is Theorem~D; the full Chern character
  $\mathrm{ch}(\mathcal{V}_g)$
  recovers the spectral discriminant~$\Delta_{\cA}$
  \textup{(}Remark~\textup{\ref{rem:spectral-characteristic-programme})}.
\item Under Verdier duality:
  $\mathsf{D}^{\mathrm{co}}(\barB^{(g)}(\cA^!)\text{-}\mathsf{comod})
  \simeq
  \mathsf{D}^{\mathrm{co}}(\barB^{(g)}(\cA)\text{-}\mathsf{comod})^{\mathrm{op}}$.
\end{enumerate}
This is the \emph{categorical lift} of
Theorems~A--D to all genera: Theorem~B
gives the equivalence, Theorem~D gives its first Chern class,
and the full equivalence is the modular functor of~$\cA$.
\end{conjecture}

\subsection{The discriminant as spectral invariant}
\index{discriminant!spectral invariant}

The shared discriminant $\Delta(x) = (1 - 3x)(1+x)$ across
$\mathfrak{sl}_2$, Virasoro, and $\beta\gamma$ systems
(Theorem~\ref{thm:ds-bar-gf-discriminant}) is not a coincidence
of generating functions.  It is the \emph{first genuinely
nontrivial non-scalar characteristic class} of the modular Koszul
object: where $\kappa(\cA)$ is a scalar measuring the curvature
of $d^2$, the discriminant $\Delta_{\cA}(x)$ records the branch
geometry of the dual Hilbert series.

Distinct chiral theories can lie on the same modular spectral
sheet --- sharing the same $\Delta$ --- even when their local
operator content is entirely different.  The underlying reason
is functorial: Drinfeld--Sokolov reduction changes the growth pole
(i.e.\ the leading OPE data) while preserving the discriminant
family.  This is compatible with the observation that
$\Delta_{\cA}$ depends only on the quadratic OPE data, which is
the Koszul-dual-determining datum.

\begin{theorem}[Discriminant as spectral determinant --- verified cases;
\ClaimStatusProvedHere]
\label{thm:discriminant-spectral-verified}
\index{discriminant!spectral determinant}
For $\cA \in \{\widehat{\mathfrak{sl}}_2,\, \mathrm{Vir},\,
\beta\gamma\}$, the bar-cohomology discriminant satisfies
\[
\Delta_{\cA}(x) \;=\; \det(1 - x\, T_{\mathrm{br}})
\]
where $T_{\mathrm{br}}$ is a $d \times d$ matrix whose eigenvalues
are the reciprocals of the branch points of the algebraic Koszul
dual Hilbert series $P_{\cA^!}(x)$, with $d$ the degree of
algebraicity.
In all three cases $d = 2$, the branch points are $x = 1/3$
and $x = -1$, and
$\Delta(x) = (1 - 3x)(1 + x)$.
\end{theorem}

\begin{proof}
Each of the three Hilbert series involves $\sqrt{1 - 2x - 3x^2}$
(Theorem~\ref{thm:ds-bar-gf-discriminant}):
$P_{\widehat{\mathfrak{sl}}_2}$ and $P_{\mathrm{Vir}}$ satisfy
quadratic equations with discriminant $(1+x)(1-3x)$, while
$P_{\beta\gamma}^2 = (1+x)/(1-3x)$ has the same branch locus.

The branch points of any algebraic function $P(x)$ satisfying
$a(x)P^2 + b(x)P + c(x) = 0$ are the zeros of
$\mathrm{Disc}_P = b(x)^2 - 4a(x)c(x)$.  For
$\widehat{\mathfrak{sl}}_2$, the algebraic equation is
$x(1+x)P^2 - (1+x)P + 1 = 0$
(Theorem~\ref{thm:ds-bar-gf-discriminant}), giving
\[
\mathrm{Disc}_P = (1+x)^2 - 4x(1+x) = (1+x)(1-3x).
\]
The branch points are $x_1 = -1$ and $x_2 = 1/3$.
Setting $\alpha_i = 1/x_i$, we obtain $\alpha_1 = -1$ and
$\alpha_2 = 3$, and
\[
\det(1 - x\,T_{\mathrm{br}}) = (1 - \alpha_1 x)(1 - \alpha_2 x)
= (1+x)(1-3x) = \Delta(x).
\]
Since $P_{\mathrm{Vir}}$ and $P_{\beta\gamma}$ involve the same
radical $\sqrt{1-2x-3x^2}$, they share the same branch locus,
hence the same~$\Delta$.
\end{proof}

\begin{theorem}[Spectral discriminant --- general case;
\ClaimStatusProvedHere]\label{thm:discriminant-spectral}
\index{discriminant!transfer matrix}
\index{discriminant!spectral determinant!general case}
Let $\cA$ be a Koszul chiral algebra whose Koszul dual Hilbert
series $P_{\cA^!}(x) = \sum_{n \geq 1} \dim(\cA^!)_n\, x^n$
is either:
\begin{enumerate}[label=\textup{(\alph*)}]
\item \emph{Rational}: $P_{\cA^!}(x) = N(x)/D(x)$ with
$\gcd(N,D) = 1$, $D(0) = 1$, and $\deg D = r$; or
\item \emph{Algebraic of degree $d \geq 2$}: satisfying an
irreducible $f(x,y) \in \mathbb{Q}[x,y]$, $\deg_y f = d$,
with discriminant $\Delta(x) = \operatorname{Disc}_y(f)$
normalized so that $\Delta(0) = 1$.
\end{enumerate}
Define the \emph{transfer matrix} $T_{\mathrm{rec}}$ as follows.

\emph{Rational case.}
The coefficients $a_n = \dim(\cA^!)_n$ satisfy a
constant-coefficient linear recurrence of order~$r$
(the degree of the denominator).  Let $T_{\mathrm{rec}}$
be its companion matrix.  Then $\Delta_{\cA}(x) := D(x)$ and
\begin{equation}\label{eq:discriminant-spectral-rational}
\Delta_{\cA}(x) \;=\; \det\!\left(1 - x\, T_{\mathrm{rec}}\right).
\end{equation}

\emph{Algebraic case.}
The coefficients $a_n$ satisfy a $P$-recursive recurrence
of order~$d$:
\begin{equation}\label{eq:p-recursive-recurrence}
q_d(n)\, a_n + q_{d-1}(n)\, a_{n-1} + \cdots + q_0(n)\, a_{n-d} = 0
\end{equation}
where $q_i(n) \in \mathbb{Q}[n]$.  Let $c_i = \lim_{n \to \infty}
q_i(n)/n^{\deg q_d}$ be the leading asymptotic coefficient of~$q_i$
\textup{(}so $c_d = 1$ after normalization\textup{)}.
The asymptotic companion matrix is
\[
T_{\mathrm{rec}} =
\begin{pmatrix}
-c_{d-1} & -c_{d-2} & \cdots & -c_0 \\
1 & 0 & \cdots & 0 \\
\vdots & & \ddots & \vdots \\
0 & \cdots & 1 & 0
\end{pmatrix}
\in \mathrm{Mat}_{d \times d}(\mathbb{Q}).
\]
Then the bar-cohomology discriminant satisfies
\begin{equation}\label{eq:discriminant-spectral-general}
\Delta(x) \;=\; \det\!\left(1 - x\, T_{\mathrm{rec}}\right).
\end{equation}
In both cases, the eigenvalues of $T_{\mathrm{rec}}$ are the
reciprocals of the singularities \textup{(}poles or branch
points\textup{)} of~$P_{\cA^!}$, and the spectral
radius of $T_{\mathrm{rec}}$ equals $\dim\mathfrak{g}$ for
Kac--Moody algebras
\textup{(}Theorem~\textup{\ref{thm:dominant-branch-point})}.
\end{theorem}

\begin{proof}
\emph{Rational case.}
Write $D(x) = \prod_{i=1}^{r}(1 - \alpha_i x)$ with
$\alpha_i = 1/x_i$ the reciprocal poles.
The partial-fraction decomposition of
$P = N/D$ shows that $a_n = [x^n]P$ satisfies a
constant-coefficient linear recurrence of order~$r$
with characteristic polynomial
$\chi(t) = \prod(t - \alpha_i)$.
The companion matrix $T_{\mathrm{rec}}$ has $\chi$ as
its characteristic polynomial, so
$\det(1 - x\,T_{\mathrm{rec}}) = \prod(1 - \alpha_i x)
= D(x) = \Delta_{\cA}(x)$.

\emph{Algebraic case.}
The argument combines three classical results from analytic
combinatorics.

\emph{Step~1: Algebraic $\Rightarrow$ $P$-recursive.}
An algebraic power series is $D$-finite
(\cite{FlajoletSedgewick}, Proposition~VI.1), hence its
coefficients satisfy a $P$-recursive (holonomic) recurrence
of order~$d$, giving~\eqref{eq:p-recursive-recurrence}.

\emph{Step~2: Singularities of the generating function.}
Since $P(x)$ is algebraic of degree~$d$ with minimal polynomial
$f(x,y) = 0$, its singularities as an analytic function are
the branch points: the roots $x_1, \ldots, x_m$ of
$\Delta(x) = \operatorname{Disc}_y(f)$
(\cite{FlajoletSedgewick}, Theorem~VII.8).

\emph{Step~3: Poincar\'e--Perron theorem.}
For the $P$-recursive recurrence~\eqref{eq:p-recursive-recurrence}
with asymptotic companion matrix~$T_{\mathrm{rec}}$, the
Poincar\'e--Perron theorem
(\cite{FlajoletSedgewick}, Theorem~VIII.4)
identifies the eigenvalues of $T_{\mathrm{rec}}$ with the
asymptotic growth rates of the solutions: provided
$c_0 \neq 0$ and the eigenvalues $\alpha_1, \ldots, \alpha_d$
have distinct absolute values, the fundamental system of
solutions satisfies
$a_n^{(i)}/a_{n-1}^{(i)} \to \alpha_i$ as $n \to \infty$.

\emph{Step~4: Growth rates = reciprocals of branch points.}
By the transfer theorem for algebraic singularities
(\cite{FlajoletSedgewick}, Theorem~VI.3), the contribution
of a simple branch point~$x_j$ to the coefficients is
$[x^n]\, P(x) \sim C_j \cdot x_j^{-n} \cdot n^{-3/2}$.
The growth rates are therefore $\alpha_i = 1/x_i$, the
reciprocals of the branch points.

\emph{Synthesis.}
Combining Steps~3 and~4, the eigenvalues of $T_{\mathrm{rec}}$
are $\alpha_i = 1/x_i$.  Since $x_1, \ldots, x_m$ are the
roots of~$\Delta$ (Step~2) and $\Delta(0) = 1$:
\[
\det(1 - x\, T_{\mathrm{rec}})
= \prod_{i=1}^d (1 - \alpha_i x)
= \prod_{i=1}^d \Bigl(1 - \frac{x}{x_i}\Bigr)
= \Delta(x). \qedhere
\]
\end{proof}

\begin{remark}[Scope and verification]%
\label{rem:discriminant-spectral-scope}
\index{discriminant!Poincar\'e--Perron}
The rational case is unconditional (given rationality of~$P$);
the algebraic case additionally requires the Poincar\'e--Perron
non-degeneracy condition ($c_0 \neq 0$, eigenvalues of distinct
absolute value).  Both conditions hold for all bar-cohomology
sequences computed in this work.  The residual conjecture is
that $P_{\cA^!}$ is algebraic or rational for all Koszul chiral
algebras associated to simple~$\mathfrak{g}$.

\emph{Algebraic verification $($degree~$2)$.}
For $\widehat{\mathfrak{sl}}_2$, Virasoro, and $\beta\gamma$
(Theorem~\ref{thm:discriminant-spectral-verified}), $P$ is
algebraic of degree~$2$.
The Riordan recurrence
$(n{+}1) a_n = (2n{-}2) a_{n-1} + (3n{-}6) a_{n-2}$
gives leading coefficients $c_2 = 1$, $c_1 = -2$, $c_0 = -3$,
and $T_{\mathrm{rec}} = \bigl(\begin{smallmatrix} 2 & 3 \\ 1 &
0 \end{smallmatrix}\bigr)$ with eigenvalues $3$ and $-1$.
Check: $\det(1 - x\,T_{\mathrm{rec}}) = 1 - 2x - 3x^2
= (1-3x)(1+x) = \Delta(x)$.  The eigenvalue
$3 = \dim(\mathfrak{sl}_2)$ is the exponential growth rate
(Theorem~\ref{thm:dominant-branch-point});
$-1$ reflects the alternating correction from the central
extension.

\emph{Rational prediction for $\widehat{\mathfrak{sl}}_3$.}
Conjecture~\ref{conj:sl3-bar-gf} predicts that
$P_{\widehat{\mathfrak{sl}}_3}$ is \emph{rational} with
denominator $(1-8x)(1-3x-x^2)$, so the theorem applies
in its rational form with $r = 3$ poles.
The companion matrix of the recurrence~\eqref{eq:sl3-bar-recurrence}
has characteristic polynomial $(t-8)(t^2-3t-1)$, giving
eigenvalues $8$, $(3+\sqrt{13})/2$, $(3-\sqrt{13})/2$.
Check: $\det(1 - x\,T_{\mathrm{rec}})
= (1-8x)(1-3x-x^2) = \Delta_{\mathfrak{sl}_3}(x)$.
The spectral radius is $8 = \dim(\mathfrak{sl}_3)$,
consistent with Theorem~\ref{thm:dominant-branch-point}.
The remaining eigenvalues
$(3 \pm \sqrt{13})/2$ are the roots of the DS discriminant
factor $t^2 - 3t - 1$, shared with the $\mathcal{W}_3$ generating
function (Remark~\ref{rem:two-discriminant-families}).

\emph{Why not the Kodaira--Spencer operator on $H^1$?}
A natural first guess is
$\Delta(x) = \det(1 - x\, T_{\mathrm{KS}} \mid H^1)$
where $T_{\mathrm{KS}}$ acts on
$H^1(\bar{B}(\cA)) \cong \mathfrak{sl}_2$ ($3$-dimensional).
But Schur's lemma forces $T_{\mathrm{KS}}$ to be scalar on the
irreducible adjoint, giving $\det = (1-cx)^3$, degree~$3 \neq 2$.
The correct object is $T_{\mathrm{rec}}$ (dimension~$d$, the
degree of algebraicity), not~$H^1$ (dimension~$\dim\mathfrak{g}$).
Identifying a natural \emph{subquotient} $H^{\mathrm{red}}_1
\subset H^1$ on which a geometric operator recovers
$T_{\mathrm{rec}}$ remains open
(Conjecture~\ref{conj:discriminant-ks-operator}).

The discriminant's role as first non-scalar characteristic class
of the modular characteristic hierarchy parallels $\kappa(\cA)$ as the first
scalar characteristic number of the universal Maurer--Cartan
class $\Theta_{\cA}$ (Theorem~\ref{thm:universal-MC}).
\end{remark}

\begin{definition}[Spectral branch object]
\label{def:spectral-branch-object}
\index{spectral branch object|textbf}
\index{discriminant!branch transport operator}
A \emph{spectral branch object} for a modular Koszul chiral
algebra~$\cA$ is a finite-rank perfect complex~$V_{\cA}^{\mathrm{br}}$
together with a \emph{branch transport operator}
$T_{\mathrm{br},\cA} \colon V_{\cA}^{\mathrm{br}} \to
V_{\cA}^{\mathrm{br}}$ such that
\begin{equation}\label{eq:discriminant-det}
\Delta_{\cA}(x)
\;=\;
\det\!\bigl(1 - x\, T_{\mathrm{br},\cA}\bigr).
\end{equation}
The pair $(V_{\cA}^{\mathrm{br}}, T_{\mathrm{br},\cA})$ satisfies:
\begin{enumerate}[label=\textup{(\roman*)}]
\item $\Delta_{\cA^!} = \Delta_{\cA}$
  (the spectral discriminant is self-dual);
\item exact factorization functors preserving quadratic OPE data
  preserve $T_{\mathrm{br},\cA}$;
\item $\operatorname{rank}(V_{\cA}^{\mathrm{br}})$
  equals the degree of algebraicity of the Koszul-dual
  Hilbert series.
\end{enumerate}

\smallskip\noindent
\emph{Status.}
In the verified cases (Theorem~\ref{thm:discriminant-spectral-verified}),
$V_{\cA}^{\mathrm{br}}$ is constructed as the span of the
companion-matrix eigenvectors of $T_{\mathrm{rec}}$,
and~\eqref{eq:discriminant-det} is proved.  The general
construction of $V_{\cA}^{\mathrm{br}}$ from the bar complex
(rather than from the recurrence) remains open; see
Conjecture~\ref{conj:discriminant-ks-operator} for the
conjectural identification with a subquotient of~$H^1$.
(Part of the MC2 theorem package; Theorem~\ref{thm:mc2-full-resolution}.)
\end{definition}

\subsection{The index theorem for genus expansions}

The appearance of the $\hat{A}$-genus in the Heisenberg genus
expansion (Theorem~\ref{thm:universal-generating-function})
and the Bernoulli asymptotics of $F_g(\cA)$ are not coincidental:
they are the output of a Grothendieck--Riemann--Roch computation
on the universal curve.

\begin{theorem}[Family index theorem for genus expansions;
\ClaimStatusProvedHere]\label{thm:family-index}
\index{index theorem!modular deformation|textbf}
\index{Grothendieck--Riemann--Roch!genus expansion|textbf}
Let $\cA$ be a Koszul chiral algebra with obstruction coefficient
$\kappa(\cA)$, let $\pi \colon \overline{\mathcal{C}}_g \to
\overline{\mathcal{M}}_g$ be the universal curve with relative
dualizing sheaf~$\omega_\pi$, and let
$\mathbb{E} = R^0\pi_*\omega_\pi$ be the Hodge bundle.
Define the \emph{modular deformation complex} of~$\cA$ at genus~$g$ by
\begin{equation}\label{eq:modular-deformation-complex-concordance}
\mathcal{D}_{\cA}^{(g)}
\;:=\;
\kappa(\cA) \cdot \mathbb{E}
\;\in\; K_0(\overline{\mathcal{M}}_g) \otimes_{\mathbb{Z}} \mathbb{Q},
\end{equation}
the $K$-theory class given by $\kappa(\cA)$ copies of the
Hodge bundle.  Then:
\begin{enumerate}[label=\textup{(\roman*)}]
\item \emph{GRR formula.}
The genus-$g$ free energy is the degree-$g$ component of the
Grothendieck--Riemann--Roch pushforward:
\begin{equation}\label{eq:family-index-grr}
F_g(\cA)
\;=\;
\kappa(\cA) \cdot
\int_{\overline{\mathcal{M}}_{g,1}} \psi^{2g-2}\, c_g(\mathbb{E})
\;=\;
\kappa(\cA) \cdot
\left[
  \pi_*\!\left(
    \operatorname{ch}(\omega_\pi) \cdot
    \operatorname{Td}(T_\pi)
  \right)
\right]_g^{\mathrm{eval}},
\end{equation}
where $[\cdot]_g^{\mathrm{eval}}$ denotes evaluation of the degree-$g$ Chern character
$\operatorname{ch}_g(\mathbb{E})$ against $\psi^{2g-2}$ on
$\overline{\mathcal{M}}_{g,1}$.

\item \emph{$\hat{A}$-genus identification.}
The total generating function is
\begin{equation}\label{eq:family-index-generating}
\sum_{g \geq 1} F_g(\cA)\, x^{2g}
\;=\;
\kappa(\cA) \cdot \bigl(\hat{A}(ix) - 1\bigr)
\;=\;
\kappa(\cA) \cdot
\left(\frac{x/2}{\sin(x/2)} - 1\right),
\end{equation}
where the function $(x/2)/\sin(x/2)$ arises from
$\operatorname{Td}(T_\pi) = c_1(\omega_\pi)/(1 - e^{-c_1(\omega_\pi)})$
in the GRR formula, Wick-rotated by the positivity of
tautological intersection numbers on
$\overline{\mathcal{M}}_g$.

\item \emph{Verdier compatibility.}
Under Koszul duality $\cA \leftrightarrow \cA^!$,
$\mathcal{D}_{\cA^!}^{(g)} = \kappa(\cA^!) \cdot \mathbb{E}$,
and the complementarity
$\kappa(\cA) + \kappa(\cA^!) = \kappa_{\mathrm{total}}$
\textup{(}Theorem~\textup{\ref{thm:genus-universality}(ii))}
gives
\[
F_g(\cA) + F_g(\cA^!)
\;=\;
\kappa_{\mathrm{total}} \cdot \lambda_g^{\mathrm{FP}},
\]
the ``index-theoretic complementarity'': the total genus expansion
of a Koszul pair is a fixed multiple of the
Faber--Pandharipande tautological integral.
\end{enumerate}
\end{theorem}

\begin{proof}
The proof assembles four ingredients, each established
independently.

\emph{Step~1: Construction of $\mathcal{D}_{\cA}^{(g)}$.}
The genus universality theorem
(Theorem~\ref{thm:genus-universality}) establishes that the
genus-$g$ obstruction class factors as
$\mathrm{obs}_g(\cA) = \kappa(\cA) \cdot \lambda_g$ in
$H^{2g}(\overline{\mathcal{M}}_g)$, where
$\lambda_g = c_g(\mathbb{E})$ is the top Chern class of the
Hodge bundle.  This factorization determines
$\mathcal{D}_{\cA}^{(g)}$ uniquely in
$K_0(\overline{\mathcal{M}}_g) \otimes \mathbb{Q}$:
the obstruction class $\mathrm{obs}_g$ is the top Chern class of
$\mathcal{D}_{\cA}^{(g)}$, and the universality theorem
shows that $c_g(\mathcal{D}_{\cA}^{(g)}) =
\kappa(\cA) \cdot c_g(\mathbb{E})$, i.e.,
$\mathcal{D}_{\cA}^{(g)} = \kappa(\cA) \cdot \mathbb{E}$
as $K$-theory classes.

The construction is \emph{functorial}: $\kappa(\cA)$ is
determined by the genus-$1$ curvature of the bar complex
(a canonical invariant of~$\cA$), and $\mathbb{E}$ is determined
by the universal curve (independent of~$\cA$).  The product
$\kappa(\cA) \cdot \mathbb{E}$ is therefore a functor from
Koszul chiral algebras to $K$-theory classes on
$\overline{\mathcal{M}}_g$.

\emph{Step~2: Mumford's GRR computation.}
By Grothendieck--Riemann--Roch applied to the universal curve
$\pi \colon \overline{\mathcal{C}}_g \to \overline{\mathcal{M}}_g$
and the relative dualizing sheaf~$\omega_\pi$:
\begin{equation}\label{eq:mumford-grr}
\operatorname{ch}(\mathbb{E})
\;=\;
\operatorname{ch}(R^0\pi_*\omega_\pi)
\;=\;
1 + \pi_*\!\left(
  \operatorname{ch}(\omega_\pi) \cdot
  \operatorname{Td}(T_\pi)
\right),
\end{equation}

Serre duality on the fibers gives
$R^1\pi_*\omega_\pi \cong (R^0\pi_*\mathcal{O})^\vee \cong \mathcal{O}$,
so $\operatorname{ch}(R\pi_*\omega_\pi) = \operatorname{ch}(\mathbb{E}) - 1$.

Write $K = c_1(\omega_\pi)$ for the relative canonical class.  Since
$c_1(T_\pi) = -K$, the Todd class of the relative tangent bundle is
$\operatorname{Td}(T_\pi) = (-K)/(1 - e^K) = K/(e^K - 1)$, and the
pushforward becomes
\begin{equation}\label{eq:mumford-chern-character}
\operatorname{ch}_n(\mathbb{E})
\;=\;
\pi_*\!\left[
  \frac{K \cdot e^K}{e^K - 1}
\right]_{\text{deg } 2n+1}
\;=\;
\frac{B_{2n}}{(2n)!}\,\kappa_{2n-1}
\;+\; \delta_n
\qquad (n \geq 1),
\end{equation}
where $\kappa_{2n-1} = \pi_*(K^{2n})$ are the
Mumford--Morita--Miller classes \cite{Mumford83}
and $\delta_n$ is supported on the boundary
$\partial\overline{\mathcal{M}}_g$.
Indeed,
\[
\frac{K e^K}{e^K-1}
= \frac{K}{1 - e^{-K}}
= \operatorname{Td}(\omega_\pi)
= 1 + \frac{K}{2} + \sum_{n \geq 1} \frac{B_{2n}}{(2n)!} K^{2n},
\]
and $\pi_*(K^{2n}) = \kappa_{2n-1}$ by definition of the tautological
classes.

\emph{Step~3: Faber--Pandharipande evaluation.}
The Faber--Pandharipande formula
(Theorem~\ref{thm:mumford-formula}) evaluates the
$\psi$-$\lambda_g$ integral:
\[
\int_{\overline{\mathcal{M}}_{g,1}} \psi^{2g-2}\,\lambda_g
\;=\;
\frac{2^{2g-1}-1}{2^{2g-1}} \cdot
\frac{|B_{2g}|}{(2g)!}.
\]
This is the content of~\eqref{eq:family-index-grr}: the genus
universality theorem reduces $F_g$ to $\kappa$ times this
integral, and the integral evaluates to the Bernoulli expression
via the tautological intersection theory initiated by
Mumford~\cite{Mumford83} and completed by
Faber--Pandharipande~\cite{FP00}.

\emph{Step~4: Generating function from the Todd class.}
The Bernoulli numbers $B_{2g}$ are the Taylor coefficients of
the Todd generating function
$x/(1 - e^{-x}) = 1 + x/2 + \sum_{g \geq 1}
\frac{B_{2g}}{(2g)!}\, x^{2g}$.
In Step~2, the Todd class $\operatorname{Td}(\omega_\pi)$
produced these Bernoulli numbers as the Chern character
$\operatorname{ch}_g(\mathbb{E})$, i.e., the genus expansion
\emph{is} the GRR pushforward.
Concretely:
\begin{align}
\sum_{g \geq 1} \lambda_g^{\mathrm{FP}}\, x^{2g}
&\;=\;
\sum_{g \geq 1}
\frac{2^{2g-1}-1}{2^{2g-1}} \cdot
\frac{|B_{2g}|}{(2g)!}\, x^{2g}
\notag \\
&\;=\;
\frac{x/2}{\sin(x/2)} - 1
\;=\; \hat{A}(ix) - 1,
\label{eq:bernoulli-a-hat}
\end{align}
where the second equality is the classical identity
$t/\sin t = \sum_{n \geq 0}
\frac{(2^{2n}-2)\,|B_{2n}|}{(2n)!}\,t^{2n}$
evaluated at $t = x/2$ (equivalently, the even part of
the Wick-rotated Todd class).
Multiplying by $\kappa(\cA)$ gives~(ii).

The positivity of all coefficients
($(x/2)/\sin(x/2)$ has all positive Taylor coefficients,
unlike $\hat{A}(x)$ which alternates) reflects the positivity
of $\lambda_g^{\mathrm{FP}} > 0$ for all~$g$, which in turn
is a consequence of the Hodge index theorem on
$\overline{\mathcal{M}}_g$.

\emph{Step~5: Verdier compatibility.}
Part~(iii) follows immediately from
$\mathcal{D}_{\cA^!}^{(g)} = \kappa(\cA^!) \cdot \mathbb{E}$
and the linearity of $F_g$ in~$\kappa$.
\end{proof}

\begin{remark}[The punchline]\label{rem:family-index-punchline}
The content of the theorem is not the generating function identity
(which was already proved as
Theorem~\ref{thm:universal-generating-function}),
but its \emph{origin}: the genus expansion of a Koszul chiral
algebra is the Grothendieck--Riemann--Roch pushforward of the
Todd class of the universal curve, scaled by the bar-complex
curvature.  The $\hat{A}$-genus appears not by accident but because
$\hat{A}$ is the multiplicative genus associated to the formal
group law $F(x,y) = x + y$ --- the trivial (additive) formal
group --- and the bar-complex genus expansion is additive in
$\kappa$.  Equivalently, the Bernoulli numbers governing both the
Todd class and the $\psi$-$\lambda_g$ integrals are the same
Bernoulli numbers because both are computed by the same GRR
calculation on the universal curve.

In physical terms: the genus expansion of a chiral algebra is
an index.  The ``partition function'' $\sum F_g \hbar^{2g-2}$
is the family index of the modular deformation complex on
$\overline{\mathcal{M}}_g$, and anomaly cancellation
($\kappa_{\mathrm{tot}} = 0$) is the vanishing of this index.

The explicit universal class
$\Theta_{\cA}^{\min} = \kappa \cdot \eta \otimes \Lambda$
(Theorem~\ref{thm:explicit-theta}) gives this a precise
homotopy-algebraic meaning: $\hat{A}$ is the character of
$\Theta_{\cA}$ --- it computes how the universal MC element
distributes its weight across genera.  The Bernoulli numbers
appear because they are the Taylor coefficients of
$\operatorname{tr}(\Theta)$, not because of any formal group
coincidence.  The formal group is additive because $\kappa$
itself is additive under tensor product
(Corollary~\ref{cor:kappa-additivity}), which in turn follows
from $\Theta_{\cA \otimes \cB} = \Theta_{\cA} + \Theta_{\cB}$
in the factorization algebra sense.
\end{remark}

\begin{remark}[Scope]\label{rem:family-index-scope}
Steps~1--3 of the proof use only results established
in this monograph (genus universality,
Theorem~\ref{thm:genus-universality}) and imported results
(Mumford's GRR formula \cite{Mumford83}; Faber--Pandharipande
\cite{FP00}).  Step~4 is classical generating function analysis.
The identification $\mathcal{D}_{\cA}^{(g)} =
\kappa(\cA) \cdot \mathbb{E}$ is canonical given the genus
universality theorem; the functoriality of the construction
follows from the functoriality of $\kappa$ (determined by
genus-$1$ curvature) and $\mathbb{E}$ (determined by the
universal curve).

\emph{What remains.}
The theorem identifies $\mathcal{D}_{\cA}^{(g)}$ via its top Chern
class (the obstruction).  A stronger result would construct
$\mathcal{D}_{\cA}^{(g)}$ as a \emph{sheaf} on
$\overline{\mathcal{M}}_g$ (not just a $K$-theory class) and
identify its full Chern character with the genus expansion at all
intermediate degrees, not just the top.  This would require the
coderived Ran-space formalism of
\S\ref{sec:modular-koszul-programme} and is the subject of
Theorem~\ref{thm:universal-MC}.
\end{remark}

\begin{remark}[The GRR theorem as horizon]
\label{rem:grr-horizon}
\index{Grothendieck--Riemann--Roch!as horizon theorem}
The theorem toward which the entire development converges is
Theorem~\ref{thm:family-index}: the genus expansion of a
Koszul chiral algebra is a GRR pushforward.
The master diagram~\eqref{eq:master-square} is the engine;
the modular characteristic hierarchy is the protagonist:
the scalar genus tower and spectral invariants are proved,
while the full quintuple
$\mathcal{C}_{\cA} = (\Theta_{\cA}, \kappa(\cA), \Delta_{\cA},
\Pi_{\cA}, \mathcal{H}_{\cA})$ now has its H-level component
$\Theta_{\cA}$ \emph{proved}
\textup{(}Theorem~\textup{\ref{thm:mc2-full-resolution}},
resolving
Theorem~\textup{\ref{thm:universal-MC}}\textup{)};
and the family index theorem is the punchline.
What remains beyond this monograph ---
$C_{\mathrm{mod}}$ (shifted-symplectic complementarity),
derived Drinfeld--Kohno, the elliptic extension ---
are further consequences of the same master diagram in richer
geometric settings.
\end{remark}

\subsection{The spectral discriminant as characteristic class}
\index{discriminant!characteristic class}

The discriminant $\Delta_{\cA}(x) = \det(1 - x\,T)$ deserves
promotion from a shared generating-function curiosity to
\emph{the first genuinely nontrivial non-scalar characteristic class
of the modular Koszul object}.  Where the scalar invariant
$\kappa(\cA)$ measures the curvature of~$d^2$, the polynomial
$\Delta_{\cA}$ records the full branch geometry of the dual
Hilbert series (Theorem~\ref{thm:ds-bar-gf-discriminant}).

Distinct chiral theories can lie on the same modular spectral
sheet --- sharing the same discriminant --- even when their local
operator content differs entirely.  Drinfeld--Sokolov reduction
changes the growth pole (the leading OPE data) while preserving
the discriminant family; this is a functorial consequence of the
modular structure, not a numerical accident.
The $K$-theoretic framework that organizes the scalar
invariant~$\kappa$ (L1), the spectral discriminant~$\Delta$ (L2),
and the periodicity profile~$\Pi$ (L3) into a hierarchy of
characteristic classes of the virtual bar
family~$\mathcal{V}_{\cA} \in K_0(\overline{\mathcal{M}}_g)$
is developed in Remark~\ref{rem:structural-saturation}; the
intrinsic construction of the L2 operator from Casimir actions
on bar cohomology is Conjecture~\ref{conj:discriminant-ks-operator}.

\subsection{Chain-level / evaluation-locus derived Drinfeld--Kohno and the Bridge Theorem}

The Yangian theory (Chapter~\ref{chap:yangians}) establishes bar-cobar
recovery for the $\Eone$-chiral Yangian and the $q \mapsto q^{-1}$
transformation of the quantum parameter.  The chain-level derived
Drinfeld--Kohno theorem is now proved
(Theorems~\ref{thm:derived-dk-affine}
and~\ref{thm:derived-dk-yangian}): the $\Eone$-chiral Koszul duality
functor, composed with the Kazhdan--Lusztig (resp.\ Kazhdan)
equivalence, yields a derived equivalence with braiding reversal
$R_q \mapsto R_{q^{-1}}$ for both affine algebras and Yangians.
The factorization-level statement on the evaluation locus is also
proved (Theorem~\ref{thm:factorization-dk-eval}).

\medskip\noindent
\textbf{The Bridge Theorem.}
The full factorization-categorical extension is the central remaining
target of the modular Koszul duality project
(Conjecture~\ref{conj:full-dk-bridge}, promoted to
Remark~\ref{rem:bridge-theorem-programme}).
The Bridge Criterion Theorem
(Theorem~\ref{thm:bridge-criterion}) shows that the
full triple equivalence
\[
\Factord(X;\cA)
\;\simeq\;
\operatorname{Mod}^{\mathrm{comp}}(\Ydg_{\cA})
\;\simeq\;
\operatorname{Rep}^{\mathrm{spec}}\!
  \bigl(\QGspec(R_{\cA})\bigr)^{\mathrm{op}}
\]
follows from three inputs:
\emph{full $\mathcal{O}$-Koszulness}~(B1),
\emph{Yangian tower completion}~(B2), and
\emph{spectral quantum-group comparison}~(B4).
An independent path through \emph{boundary/line realization}~(B3)
via the HT literature
(Remark~\ref{rem:boundary-line-bridge-path})
provides a potentially more constructive route to B4.
The remaining work is staged rather than monolithic:

\begin{conjecture}[Factorization-categorical DK/KL bridge (DK-3);
  \ClaimStatusConjectured]
\label{conj:derived-drinfeld-kohno}
\index{Drinfeld--Kohno!derived}
There exists an equivalence of $\Eone$-factorization categories
\[
\operatorname{Fact}_{\Eone}(\mathcal{Y}(\mathfrak{g}))
\;\simeq\;
\operatorname{Fact}_{\Eone}(U_q(\mathfrak{g}))^{\mathrm{op}}
\]
extending the chain-level equivalence of
Theorem~\textup{\ref{thm:derived-dk-yangian}}
from evaluation modules to all objects of Yangian
category~$\mathcal{O}$, and identifying the Yangian
$\Eone$-factorization structure with the factorisation quantum
group structure of Latyntsev~\cite{Latyntsev23}.
\end{conjecture}

\begin{remark}[Scope and ladder]\label{rem:derived-dk-scope}
The Yangian extension problem is best read as a DK ladder:
\begin{enumerate}[label=\textup{(\alph*)}]
\item \textbf{DK-0 (proved):} braided monoidal bar complexes,
  Verdier $= R$-matrix inversion, and the chain-level DK square
  (Theorem~\ref{thm:monoidal-module-koszul};
  Proposition~\ref{prop:yangian-module-koszul}(iii);
  Theorem~\ref{thm:derived-dk-yangian}).
\item \textbf{DK-1 (proved):} the factorization-level statement on the
  evaluation locus (Theorem~\ref{thm:factorization-dk-eval}).
\item \textbf{DK-2/3 (proved, evaluation-generated core, all simple types, domain: $\mathcal{O}_{\mathrm{poly}}$):}
  The domain is the polynomial subcategory
  $\mathcal{O}_{\mathrm{poly}}$
  (Corollary~\ref{cor:dk-poly-catO}), not the full
  Category~$\mathcal{O}$.
  Two independent proofs.  Both rely on Yangian Koszulness
  (Proposition~\ref{prop:yangian-koszul}), which uses Molev's PBW
  theorem \cite{molev-yangians} as external input.
  \emph{Type~$A$ thick generation}:
  Theorem~\ref{thm:dk-fd-typeA} resolves DK-2/3 for
  $\operatorname{Rep}_{\mathrm{fd}}$ of $Y(\mathfrak{sl}_N)$
  via thick generation
  (Proposition~\ref{prop:dk2-thick-generation-typeA}).
  \emph{Sectorwise convergence (all types)}:
  Corollary~\ref{cor:dk23-all-types} resolves DK-2/3
  for the evaluation-generated core at all
  simple types, via the $\infty$-categorical sectorwise
  convergence theorem
  (Theorem~\ref{thm:h-level-factorization-kd}), bypassing
  thick generation.  The remaining gap to
  Conjecture~\ref{conj:derived-drinfeld-kohno} is the
  completed/coderived extension beyond the generated core
  (Corollary~\ref{cor:dk-partial-conj}).
\item \textbf{DK-4/DK-5 (the Bridge Theorem project):} the dg-shifted Yangian and
  full triple-bridge completion
  (Conjecture~\ref{conj:full-dk-bridge}).
  Decomposed into four sub-targets~B1--B4
  (Remark~\ref{rem:bridge-theorem-programme}):
  full $\mathcal{O}$-Koszulness~(B1),
  tower completion~(B2),
  boundary/line realization~(B3),
  spectral QG comparison~(B4).
  The Bridge Criterion Theorem
  (Theorem~\ref{thm:bridge-criterion}) gives
  B1$+$B2$+$B4 $\Rightarrow$ full bridge.
  The boundary/line path~(B3) provides an independent route to B4.
\end{enumerate}
The key control point is that the evaluation-locus theorem is already
theorematic; the category-$\mathcal{O}$ extension and the dg-shifted
comparison are downstream stages, not hidden inside the same claim.
The six-lemma proof architecture decomposes the residual gap into
two proved lemmas (evaluation-locus bridge, Verdier $= R^{-1}$) and
four open kill targets (Lemmas~3--6).
See \S\ref{sec:derived-dk} for the full treatment,
Remark~\ref{rem:bridge-theorem-programme} for the full decomposition,
and Remark~\ref{rem:dk-proved-vs-remains} for the precise staged
status.
The exact extra DK-2/DK-3 packets are recorded in
Remark~\ref{rem:yangian-dk23-inputs}.

\emph{Homotopy template} (Convention~\ref{conv:hms-levels}): Type~II
(equivalence of homotopy categories).  At H-level: an equivalence of
$\Eone$-monoidal $\infty$-categories.  The M-level model is the braided
monoidal bar-cobar dg category; the S-level shadow is the classical DK
equivalence on evaluation modules.
(Contributing to Conjecture~\ref{conj:master-dk-kl}.)

DK is naturally an $\Eone$-factorization theorem: it compares two
$\Eone$-level objects (KZ monodromy and $R$-matrix data).  The step
from $\Eone$ to full $\Etwo$ coherence requires Drinfeld associator
input and is governed by the Grothendieck--Teichm\"uller group
(Remarks~\ref{rem:e1-e2-obstruction}
and~\ref{rem:grothendieck-teichmuller}).  The modular layer (genus
deformation of $\Eone$-factorization) is the content of the modular
homotopy theory project
(Definition~\ref{def:modular-homotopy-theory-intro}).
\end{remark}

\subsection{Summary}

The historical rhyme is this.  Classical operadic Koszul duality
says: quadratic presentation plus bar-cobar formalism produces
homotopy theory.  The theory sought here says: factorization
presentation on curves plus Verdier duality on modular
compactifications produces \emph{modular homotopy theory}.  Just as
Drinfeld--Kohno related KZ monodromy to quantum groups, the theory
sought here would relate configuration-space/Deligne--Mumford
monodromy to curved chiral Koszul duality.  And just as \'etale
cohomology supplied the correct ambient cohomology theory for
Frobenius and zeta functions, factorization homology on
$\operatorname{Ran}(X)$ together with Verdier duality on
$\overline{\mathcal{M}}_{g,n}$ is the correct ambient cohomology
theory for OPE, monodromy, and bar-cobar duality on curves.

\subsection{The holographic completion}
\label{subsec:concordance-holographic-completion}
\index{holographic modular Koszul datum!concordance}

The modular homotopy theory of the preceding subsections acquires a
concrete physical target via HT holography.  The holographic modular
Koszul datum
$\mathcal{H}(T) = (\cA, \cA^!, \mathcal{C}, r(z), \Theta_\cA,
\nabla^{\mathrm{hol}}_{\bullet,\bullet})$
(Definition~\ref{def:holographic-modular-koszul-datum})
is the single object that should organize HT holography as precise
algebra.  It synthesizes five literatures:
\begin{enumerate}[label=\textup{(\roman*)}]
\item the bulk/boundary chiral algebra framework of
  Costello--Dimofte--Gaiotto~\cite{CDG20};
\item the line-operator category and dg-shifted Yangian of
  Dimofte--Niu--Py~\cite{DNP25};
\item the non-planar interface/minimal-model matching of
  Gaiotto--Zeng~\cite{GZ26};
\item the 3d Poisson sigma model for PVA quantization of
  Khan--Zeng~\cite{KhanZeng25};
\item the factorisation quantum-group framework of
  Latyntsev~\cite{Latyntsev23}.
\end{enumerate}
The organizing sentence: \emph{the 2026 non-planar sphere algebra
is not the endpoint; it is the binary genus-zero shadow of the
modular Maurer--Cartan class~$\Theta_\cA$.}

\noindent\emph{What is proved.}
The full $\Theta_\cA$ and clutching law are proved
(Theorem~\ref{thm:mc2-bar-intrinsic}).  The scalar shadow
$\kappa(\cA)\lambda_g$ is proved (Theorem~\ref{thm:genus-universality}).
The shadow Postnikov tower through arity~$4$ is proved, and the
quartic resonance class with clutching law is established
(Appendix~\ref{app:nonlinear-modular-shadows}).
Beyond the conjectural framework, the collision filtration
(\S\ref{subsec:collision-filtration-recovery}) establishes four
\textbf{proved recovery theorems}:
\begin{enumerate}[label=\textup{(\alph*)}]
\item \emph{Collision residue = twisting morphism}
  (Theorem~\ref{thm:collision-residue-twisting}):
  $\operatorname{Res}^{\mathrm{coll}}_{0,2}(\Theta_\cA) =
  \pi_\cA$, the universal twisting morphism, canonically producing
  the dg-shifted Yangian $r_\cA(z)$.
\item \emph{Arnold/MC at depth~$2$ = $A_\infty$-YBE}
  (Theorem~\ref{thm:collision-depth-2-ybe}): the CYBE from the
  Arnold relation on $\overline{\mathcal{M}}_{0,4}$.
\item \emph{Shadow connection = KZ}
  (Theorem~\ref{thm:shadow-connection-kz}): for
  $\cA = \widehat{\fg}_k$, the genus-$0$ shadow connection is
  the KZ connection $\nabla^{\mathrm{KZ}}$, with flatness
  from the MC equation; BPZ recovery for Virasoro
  (Proposition~\ref{prop:shadow-connection-bpz}).
\item \emph{Quartic obstruction = $L_\infty$ obstruction}
  (Theorem~\ref{thm:quartic-obstruction-linf}):
  $[o_4(\cA)] = [\mathfrak{R}^{\mathrm{mod}}_4(\cA)]$;
  vanishing for classes~$G,L,C$; non-vanishing for class~$M$.
\end{enumerate}
The complete Heisenberg datum is computed explicitly
(Computation~\ref{comp:heisenberg-holographic-datum}).
The holographic spectral sequence
(Construction~\ref{constr:holographic-spectral-sequence})
classifies shadow depth by collapse stage.
The five shadows of $\Theta_\cA$---scalar, $r$-matrix, connection,
quartic, clutching---are five projections of one MC element
(Remark~\ref{rem:five-shadows-synthesis}).

\noindent\emph{What remains conjectural.}
The five programme-level targets:
boundary--defect realization
(Conjecture~\ref{conj:boundary-defect-realization}),
Yangian-shadow in full derived generality
(Conjecture~\ref{conj:yangian-shadow-theorem}),
sphere reconstruction beyond KZ/BPZ
(Conjecture~\ref{conj:sphere-reconstruction}),
quartic resonance obstruction as the first test beyond scalar
(Conjecture~\ref{conj:quartic-resonance-obstruction}), and
singular-fiber descent for minimal models
(Conjecture~\ref{conj:singular-fiber-descent}).
The HT deformation-quantization functor
(Conjecture~\ref{conj:ht-deformation-quantization}) and the
higher-genus shadow connection at $g \geq 2$
(Definition~\ref{def:modular-shadow-connection}) remain
conjectural.
